\chapter{概率论与数理统计}
\label{chap:probability}

从一批产品中抽取若干件，记录其中的不合格品数，可以提出两种不同的问题。若已知这批产品的构成，抽样会得到怎样的结果？若只知道抽检结果，又能对整批产品作出怎样的判断？前一个问题从模型走向可能的数据，后一个问题从数据研究未知的总体。概率与统计将在这里相接。

一次抽取的结果可以有许多种。先确定一次完整试验记录什么，再研究哪些结果满足所关心的条件，就有了计数、事件和概率的共同起点。已知一部分信息以后，还可以重新考察剩余结果的概率；独立性则说明，一种信息在什么意义下不改变另一事件的概率。

把抽到的不合格品数、等待的次数或测量结果记成数值，就得到随机变量。研究它的分布，可以知道各类结果怎样分配；研究期望和方差，可以描述平均水平与波动。取值不同的随机变量，有时遵循相同的运算关系；前面建立的函数和线性思想将在这些关系中继续发挥作用。

统计面对的是实际取得的数据。抽样方式决定数据能说明什么，图表与数字特征帮助整理观察，抽样分布则把样本之间的波动也纳入研究。大数定律说明平均与频率的稳定依据，参数估计、假设检验和回归据此提出各自的判断任务。一次判断的意义，要由它所依据的模型、数据取得方式及计算过程共同确定。

本章先建立有限模型，再进入具体的离散与连续分布。主线使用高中代数、函数及极限的基本认识；选讲调用向量空间、内积和积分，从新的层面解释同一批问题。抽检、重复试验与成对数据会在后文反复出现，每次承担不同的任务。方法随着这些任务产生，研究到依据清楚、问题得到回答时，便可以收束。

\FloatBarrier
\newpage
\section{计数原理}
\label{sec:prob-counting}

计数之前，先要确定所数的对象，以及怎样的记录算作同一个结果。选择两名同学与给两名同学安排不同职务，涉及同样的人，却有不同的结果识别方式。下面先研究组织计数的共同原则，再建立常用公式。

\subsection{分类计数}

\begin{proposition}[加法原理]
若一个有限集合\(S\)分成两两不交的子集\(S_1,\ldots,S_m\)，则
\[
 |S|=|S_1|+\cdots+|S_m|.
\]
\end{proposition}

每个对象恰在一个类别中出现，因而将各类的个数相加，便把所有对象各数一次。这也是分类计数的依据：类别需要覆盖全部对象，类别之间则不能重复。

若原有分类相互重叠，可以先扣除重复。对有限集合\(A,B\)，将并集分成\(A\setminus B\)、\(A\cap B\)、\(B\setminus A\)，得到
\begin{equation}\label{eq:prob-counting-inclusion}
 |A\cup B|=|A|+|B|-|A\cap B|.
\end{equation}
同理，三个集合满足
\begin{align*}
 |A\cup B\cup C|={}&|A|+|B|+|C|\\
 &-|A\cap B|-|A\cap C|-|B\cap C|+|A\cap B\cap C|.
\end{align*}
属于三个集合的对象，先被加三次、减三次，最后再加一次；其他对象也恰被保留一次。这称为{\heiti 容斥原理}。若所求对象的反面更易计数，也可以用总数减去反面对象数。

\begin{example}
某班有40人，参加数学社团的有18人，参加物理社团的有16人，两者都参加的有7人。求至少参加一个社团、恰参加一个社团和两个社团都未参加的人数。
\end{example}

\subsection{分步计数}

设完成一件事需要依次作出\(m\)次选择。若第一步有\(n_1\)种选法，并且对每一种已经完成的前序选择，第\(j\)步都有\(n_j\)种选法，则完整结果数为
\begin{equation}\label{eq:prob-counting-product}
 n_1n_2\cdots n_m.
\end{equation}
这称为{\heiti 乘法原理}。两步时，每个第一步结果对应\(n_2\)个完整结果，由加法原理共得\(n_1n_2\)个；一般情形由归纳法得到。

可以用树状图记录选择过程，每条从起点到终点的路径对应一个完整结果。若后一步的选法数随前面的具体选择而改变，就分别计算各分支，再把结果相加。例如，第一步选甲类或乙类对象，分别有2种、3种；选甲类后第二步有4种，选乙类后第二步有5种，则结果数为\(2\cdot4+3\cdot5\)。

\begin{figure}[!ht]
\centering
\includegraphics[width=.50\linewidth]{probability-counting-tree}
\caption{按类别合并的计数树，\(A,B\)分别表示甲、乙两类；边上数字为相应选择数。}
\label{fig:prob-counting-tree}
\end{figure}
\FloatBarrier

图\ref{fig:prob-counting-tree}先分类、再分步，类别间用加法合并，各类内部用乘法计数。

\begin{example}
用数字\(0,1,2,3,4\)组成没有重复数字的三位数。分别求全部三位数、偶数和大于300的三位数的个数。比较各问适合先确定哪个位置。
\end{example}

\subsection{排列}

\begin{definition}
从\(n\)个不同对象中取出\(r\)个，按次序排成一列，称为一个{\heiti 排列}；排列的总数记为\(A_n^r\)。取出全部对象的排列称为{\heiti 全排列}。
\end{definition}

依次选择各位置的对象，分别有\(n,n-1,\ldots,n-r+1\)种选法，所以
\[
 A_n^r=n(n-1)\cdots(n-r+1)=\frac{n!}{(n-r)!},\qquad A_n^n=n!.
\]
这里\(n!=1\cdot2\cdots n\)，并约定\(0!=1\)。若每个位置可以重复使用同一对象，长度为\(r\)的有序记录共有\(n^r\)个。

若一列中有\(n\)个元素，其中各类相同元素分别出现\(n_1,\ldots,n_k\)次，\(n_1+\cdots+n_k=n\)，则不同排列数为
\[
 \frac{n!}{n_1!\cdots n_k!}.
\]
给同类元素临时加上不同标记，共有\(n!\)个有标记排列；去掉标记后，每个原排列恰对应\(n_1!\cdots n_k!\)种标记方式，因而得到此式。

\begin{example}
六名同学排成一列，甲、乙是其中两人。分别求甲乙相邻、甲乙不相邻，以及甲在乙左侧的排列数。再求由字母\(A,A,A,B,B,C\)组成的不同字符串数。
\end{example}

\subsection{组合}

\begin{definition}
从\(n\)个不同对象中选出\(r\)个组成一个子集，称为一个{\heiti 组合}，组合数记为\(C_n^r\)或\(\binom nr\)。
\end{definition}

每个\(r\)元子集可以排成\(r!\)个排列，故
\begin{equation}\label{eq:prob-combinations}
 \binom nr=\frac{A_n^r}{r!}=\frac{n!}{r!(n-r)!}.
\end{equation}
这里除以\(r!\)，是因为每个最终子集都恰好对应\(r!\)个有序记录。通过另一种记录计数时，应先确定每个最终结果对应多少条记录；这个对应数相同，才能直接除以它。

\begin{example}
从6名同学中选3人组成小组，并指定其中一人为组长。分别按“先选小组，再选组长”和“先选组长，再选组员”计数。若先将三人排成一列，以首人为组长，再忘去后两人的次序，每个结果被重复计数几次？
\end{example}

取\(r\)个对象与留下\(n-r\)个对象相对应，所以\(\binom nr=\binom n{n-r}\)。再固定一个对象，按是否选入它分类，得到
\[
 \binom nr=\binom{n-1}{r}+\binom{n-1}{r-1}.
\]
计数中约定\(r<0\)或\(r>n\)时\(\binom nr=0\)，便可统一书写求和范围。

将对象分成大小为\(a,b\)的两组，按第一组选入的个数计数，得到{\heiti Vandermonde恒等式}
\begin{equation}\label{eq:prob-vandermonde}
 \sum_k\binom ak\binom b{r-k}=\binom{a+b}r.
\end{equation}
等式两边计算的是同一批\(r\)元子集。类似地，把\(n\)个不同对象分入有标记的各组，组大小为\(n_1,\ldots,n_k\)，有\(n!/(n_1!\cdots n_k!)\)种分法。若组不带标记，还需根据实际重复次数处理；例如把四人分成两个无标记二人组，共有\(\binom42/2\)种。

\subsection{二项式定理}

展开\((a+b)^n\)时，每个因子中选\(a\)或\(b\)。含\(k\)个\(b\)的项来自\(n\)个因子中选出\(k\)个，因而
\begin{theorem}[二项式定理]
对非负整数\(n\)，有
\begin{equation}\label{eq:prob-binomial-theorem}
 (a+b)^n=\sum_{k=0}^n\binom nk a^{n-k}b^k.
\end{equation}
\end{theorem}

按\(b\)的次数从0到\(n\)排列，第\(k+1\)项为
\[
 T_{k+1}=\binom nk a^{n-k}b^k,\qquad k=0,1,\ldots,n.
\]
若\(a,b\)本身含有变量\(x\)，则\(x\)的指数需要合并计算。例如取\(a=x^2,b=2/x\)，通项为\(\binom nk2^kx^{2n-3k}\)。求某次幂的系数时，先确定对应的\(k\)，再计算该项的数值系数，组合数只是系数中的一个因子。

取\(a=b=1\)，得到\(\sum_k\binom nk=2^n\)；取\(a=1,b=-1\)，得到\(n\ge1\)时奇数项与偶数项组合数之和相等。

还有两条后面会用到的恒等式：
\begin{equation}\label{eq:prob-binomial-moments}
 \sum_{k=0}^n k\binom nk=n2^{n-1},\qquad
 \sum_{k=0}^n k(k-1)\binom nk=n(n-1)2^{n-2}\quad(n\ge2).
\end{equation}
第一式数一个子集及其中一个被标记的元素：先选标记元素，再从其余元素中任选一个子集。第二式则标记两个有先后次序的不同元素。它们也分别对应
\(k\binom nk=n\binom{n-1}{k-1}\)与
\(k(k-1)\binom nk=n(n-1)\binom{n-2}{k-2}\)。

\begin{example}[2025年上海卷，题面整理]
求\((2x-1)^5\)展开式中\(x^3\)的系数。
\end{example}

\begin{example}
求\((x^2+2/x)^9\)展开式的常数项、\(x^6\)项的系数及各项系数之和。
\end{example}

\begin{exercise}
\mbox{}\par\nobreak

（1）用数字\(0,1,2,3,4,5\)组成没有重复数字的四位数，求其中能被5整除的个数。

（2）七人中选三人，其中甲、乙至少选入一人，有多少种选法？若三人还需分别担任三个不同职务，又有多少种？

（3）将六人分成有标记的两个三人组，再分成无标记的两个三人组，分别有多少种分法？将六人分为一个二人组和一个四人组时，是否还需要除以2？

（4）求\((x^2-2/x)^6\)的常数项，并求各项系数之和。

（5）用对子集的计数证明\(\sum_{k=0}^n\binom nk^2=\binom{2n}n\)。

（6）有三类对象，数量分别为\(a,b,c\)。证明\(\sum_{i+j+k=r}\binom ai\binom bj\binom ck=\binom{a+b+c}r\)。

（7）（2023年全国甲卷理科，题面整理）5名志愿者参加公益活动，周六、周日每天安排2人。恰有1人两天都参加，共有多少种安排？
\end{exercise}

\FloatBarrier
\newpage
\subsection{计数对象的变换*}

把\(n\)个相同物品放入\(r\)个不同盒子，各盒允许为空，等价于求
\[
 x_1+\cdots+x_r=n,\qquad x_i\in\{0,1,2,\ldots\}
\]
的解数。在\(n\)个相同标记之间及两端放入\(r-1\)块隔板，每段的标记数就是相应的\(x_i\)。反过来，由一组\(x_i\)也能唯一恢复标记与隔板的次序。因此解数为
\[
 \binom{n+r-1}{r-1}.
\]
若每盒至少一个，令\(y_i=x_i-1\)，便转化为\(\sum y_i=n-r\)，解数为\(\binom{n-1}{r-1}\)。这里盒子的标记决定坐标的次序，不能再把不同坐标排列随意视为相同。

另一个例子是圆排列。\(n\)个不同的人围圆桌就座，只把整体旋转视为同一排列。固定其中一人的位置，其余人依次排列，共有\((n-1)!\)种。镜面反向通常仍是不同座次；若题目还将反向识别为同一结果，需要另外处理这种识别。

这些计算都先建立了对象之间的一一对应，公式随后由已知计数得到。

\begin{exercise}
\mbox{}\par\nobreak

（1）把10个相同球放入4个不同盒子，求允许空盒、每盒至少一个及第一盒至少两个时的分法数。

（2）从\(r\)种口味的糖果中购买\(n\)颗，同种口味不加区分。把购买方案与非负整数解对应，求方案数。

（3）六人围圆桌就座，甲、乙不能相邻。求不同座次的个数。
\end{exercise}

\FloatBarrier
\newpage
\section{概率模型}
\label{sec:prob-models}

计数确定可能结果的个数，概率还需要说明各结果怎样分配权重。先把一次试验记录清楚，才能判断所用的计数是否适合当前概率计算。

\subsection{样本空间}

\begin{definition}
一次随机试验的一个完整可能结果称为{\heiti 样本点}，全体样本点组成的集合称为{\heiti 样本空间}，通常记为\(\Omega\)。
\end{definition}

例如，连续抛掷两次硬币，分别用\(H,T\)记录正反面，可以取
\[
 \Omega=\{HH,HT,TH,TT\}.
\]
若只关心正面次数，也可以记录为\(\{0,1,2\}\)，但两个模型保留的信息不同。前者可以区分第一次正面还是第二次正面，后者把\(HT,TH\)合成了同一种数值记录。三个数值并不会因此自动成为等可能结果。

本章先研究有限样本空间。建立它时，应使每次试验恰能归入一个样本点，所有可能结果均在其中。

\subsection{事件}

\begin{definition}
样本空间的一个子集称为{\heiti 事件}。试验结果属于这个子集时，就说该事件发生。\(\Omega\)称为必然事件，\(\varnothing\)称为不可能事件。
\end{definition}

事件\(A\cup B\)表示\(A,B\)至少一个发生，\(A\cap B\)表示两者同时发生。补事件\(A^c=\Omega\setminus A\)表示\(A\)不发生；\(A\setminus B=A\cap B^c\)表示\(A\)发生而\(B\)不发生。

若\(A\cap B=\varnothing\)，称\(A,B\)互斥。若\(A\subseteq B\)，则\(A\)发生必有\(B\)发生。例如在两次抛币中，“两次正面”包含于“至少一次正面”。事件的关系由它所含的样本点确定，可以直接使用集合运算。

例如两次抛币，事件\(A=\{HH,HT\}\)表示第一次正面，\(B=\{HH,TH\}\)表示第二次正面。“恰一次正面”便是\((A\setminus B)\cup(B\setminus A)=\{HT,TH\}\)。

若\(\Omega\)有\(N\)个样本点，事件共有\(2^N\)个：每个样本点都可以选择属于或不属于所取的子集。因此，样本点的个数与事件的个数是不同的计数对象。

\begin{example}
从编号1至5的球中依次不放回取两球。分别用有序对和无序集合表示样本空间，并判断哪种记录能表示事件“第一次取到1号球”。
\end{example}

\newpage
\subsection{概率分配}

对有限样本空间\(\Omega=\{\omega_1,\ldots,\omega_N\}\)，给各样本点分配数\(p_i\)，满足
\[
 p_i\ge0,\qquad \sum_{i=1}^N p_i=1.
\]
定义事件\(A\)的{\heiti 概率}为
\begin{equation}\label{eq:prob-finite-probability}
 P(A)=\sum_{\omega_i\in A}p_i.
\end{equation}
于是\(P(\Omega)=1\)、\(P(\varnothing)=0\)。互斥事件的并，其概率等于概率之和。由集合分解还得到
\begin{align}
 P(A^c)&=1-P(A),\label{eq:prob-complement}\\
 P(A\cup B)&=P(A)+P(B)-P(A\cap B).\label{eq:prob-addition}
\end{align}
若\(A\subseteq B\)，由\(B=A\cup(B\setminus A)\)得到\(P(A)\le P(B)\)。这些性质都来自非负权重的求和。

例如两次独立公平抛币的四个序列各分配\(1/4\)，正面次数\(0,1,2\)的概率分别为\(1/4,1/2,1/4\)。记录合并时，需要把相应样本点的概率相加。

样本空间、其事件集合及概率分配共同组成一个{\heiti 有限概率空间}。机制不同，可以在同一结果集合上给出不同的概率分配。

\subsection{古典概型}

若有限样本空间中各样本点等可能，则每个样本点的概率为\(1/N\)，式\eqref{eq:prob-finite-probability}成为
\begin{equation}\label{eq:prob-classical}
 P(A)=\frac{|A|}{|\Omega|}.
\end{equation}
这种模型称为{\heiti 古典概型}。加法、乘法原理与排列组合公式在这里提供分子、分母的计数。

例如从\(n\)个不同对象中均匀选取\(r\)个，所有\(r\)元子集等可能；依次不放回抽取时，所有长度为\(r\)的无重复序列也等可能。每个子集都对应\(r!\)个序列，因此两种记录下，同一个事件的概率相同。计算时只需前后一致地选择一种记录。

\begin{example}[2023年全国甲卷文科，题面整理]
某校文艺部有4名学生，高一、高二各2名。从中随机选2名组织文艺汇演，求两人来自不同年级的概率。
\end{example}

\begin{example}[2024年全国甲卷文科，题面整理]
甲、乙、丙、丁各出场一次，次序由随机抽签确定。求丙不是第一个出场，且甲或乙最后出场的概率。
\end{example}

\newpage
\subsection{频率与模拟}

在\(n\)次试验中，事件\(A\)发生\(N_A\)次，称\(N_A/n\)为\(A\)在这组试验中的{\heiti 频率}。它由实际记录得到；模型中的\(P(A)\)则由概率分配确定。若模型含未知参数，可以用观测频率估计它。

随机模拟可以给出新的试验记录。例如用公平随机整数\(0,1,\ldots,9\)，将\(0,1,2,3\)记为成功，便得到成功概率为\(0.4\)的试验。独立重复这个步骤，记录不同试验次数下的成功频率，就能观察波动如何变化。

\begin{figure}[!ht]
\centering
\includegraphics[width=.78\linewidth]{probability-frequency}
\caption{成功概率为\(0.4\)的独立重复试验的一条模拟记录。虚线表示模型概率。}
\label{fig:prob-frequency}
\end{figure}

模拟中的频率仍会波动。为什么较大的试验次数通常能带来更可靠的估计，将由后面的方差与大数定律说明。

\begin{exercise}
\mbox{}\par\nobreak

（1）独立抛掷两枚公平骰子。分别用有序点数对与点数和记录结果，说明哪一种记录的样本点等可能，并求点数和为7的概率。

（2）设\(P(A)=0.6\)、\(P(B)=0.5\)、\(P(A\cup B)=0.7\)，求\(P(A\cap B)\)、\(P(A\setminus B)\)和\(P(A^c\cap B^c)\)。

（3）从编号1至6的球中不放回抽取3个。分别按有序抽取和无序子集计数，求其中恰有两个偶数的概率。

（4）某次模拟做1000次试验，事件发生397次。再做1000次时，是否必须恰好发生403次，才能使总频率等于模型给定的\(0.4\)？说明模型概率与实际记录各自表示什么。
\end{exercise}

\FloatBarrier
\newpage
\section{条件概率}
\label{sec:prob-conditional}

已经知道一件事发生，相当于获得了关于试验结果的新信息。概率计算需要在这部分结果中重新分配权重。条件可能来自试验的先后过程，也可能来自对同一次结果的进一步观察。

\subsection{条件下的模型}

\begin{definition}
若\(P(B)>0\)，称
\begin{equation}\label{eq:prob-conditional}
 P(A\mid B)=\frac{P(A\cap B)}{P(B)}
\end{equation}
为已知\(B\)发生时\(A\)的{\heiti 条件概率}。
\end{definition}

在有限模型中，保留\(B\)内的样本点，将其原权重都除以\(P(B)\)，总权重便重新成为1。\(A\)在这个模型中的部分是\(A\cap B\)，因此得到上式。若原样本点等可能，还可写为\(P(A\mid B)=|A\cap B|/|B|\)，即在\(B\)内重新计数。

例如公平抛两次硬币，已知至少一次正面，剩下\(HH,HT,TH\)三个等可能结果，其中两次正面的概率为\(1/3\)。若已知第一次正面，剩下的是\(HH,HT\)，条件概率则为\(1/2\)。不同信息对应不同条件事件。

\begin{example}[2023年全国甲卷理科，题面整理]
某地中学生中，60\%爱好滑冰，50\%爱好滑雪，70\%至少爱好其中一项。随机调查一名学生，已知他爱好滑雪，求他也爱好滑冰的概率。
\end{example}

\subsection{乘法公式}

把定义移项，即得
\begin{equation}\label{eq:prob-multiplication}
 P(A\cap B)=P(B)P(A\mid B).
\end{equation}
多阶段事件可以连续使用此式。例如
\[
 P(A\cap B\cap C)=P(A)P(B\mid A)P(C\mid A\cap B),
\]
其中各条件事件均有正概率。树状图的每条完整路径，把沿途的条件概率相乘；不同终点路径互斥时，再把路径概率相加。

袋中有3个红球、2个白球，不放回抽两球，所得两球异色的概率为
\[
 \frac35\frac24+\frac25\frac34=\frac35.
\]
第二次抽取的条件概率由第一次实际拿走的颜色决定。放回抽取时，袋中构成恢复，路径上的相应分母和分子便随之改变。

\begin{example}[2022年天津卷，题面整理]
从含4张A的52张扑克牌中依次随机取两张，不放回。求两张都是A的概率，以及第一张是A时第二张也是A的条件概率。
\end{example}

\subsection{全概率公式}

\begin{theorem}[全概率公式]
设\(B_1,\ldots,B_m\)两两互斥，其并为\(\Omega\)，且\(P(B_i)>0\)，则
\begin{equation}\label{eq:prob-total}
 P(A)=\sum_{i=1}^mP(B_i)P(A\mid B_i).
\end{equation}
\end{theorem}

\begin{proof}
\(A\)是互斥事件\(A\cap B_i\)的并，先用可加性，再用乘法公式即可。
\end{proof}

下面固定一个抽检模型供本节使用。产品来自甲、乙两条生产线，来源比例分别为\(0.4,0.6\)，不合格率分别为\(0.02,0.05\)。从混合产品中随机抽一件，令\(D\)表示不合格，则
\[
 P(D)=0.4\cdot0.02+0.6\cdot0.05=0.038.
\]
总体不合格率是各来源不合格率按来源比例的加权平均。图中路径同时保留了来源与检验结果。

\begin{figure}[!ht]
\centering
\includegraphics[width=.58\linewidth]{probability-tree}
\caption{按来源分类，再记录检验结果；每条路径的概率由两个条件因子相乘得到。}
\end{figure}

\subsection{贝叶斯公式}

\begin{theorem}[贝叶斯公式]
在全概率公式的条件下，若\(P(A)>0\)，则
\begin{equation}\label{eq:prob-bayes}
 P(B_j\mid A)=
 \frac{P(B_j)P(A\mid B_j)}{\sum_iP(B_i)P(A\mid B_i)}.
\end{equation}
\end{theorem}

\begin{proof}
在条件概率定义中代入乘法公式与全概率公式即得。
\end{proof}

在上述抽检模型中，已知产品不合格，它来自甲线的概率为
\[
 P(\text{甲}\mid D)=\frac{0.4\cdot0.02}{0.038}=\frac4{19}.
\]
抽检以前，甲线占全部产品的\(0.4\)；知道不合格后，应比较各线贡献的不合格产品数，比例于是改变。分母负责把这些贡献重新归一化。

把上述比例换成每万件产品的构成，可以更直接地看到所取的分母：
\[
\begin{array}{c|rr|r}
 &\text{不合格}&\text{合格}&\text{合计}\\\hline
 \text{甲线}&80&3920&4000\\
 \text{乙线}&300&5700&6000\\\hline
 \text{合计}&380&9620&10000
\end{array}
\]
表中数值按模型比例列出。甲线有80件不合格产品，已知来自甲线时分母取4000，已知不合格时则取380。分子对应同一交事件，分母随条件改变。

\begin{example}[据2025年北京卷改编]
甲、乙两校各随机调查100名学生，分别有80人、75人答对某四选一选择题。用频率估计答对概率。甲校学生掌握知识点时必答对，乙校学生掌握时答对概率为\(0.85\)；两校未掌握者均等可能地猜选四个选项。分别估计两校学生掌握知识点的比例，并比较大小。
\end{example}

\begin{exercise}
\mbox{}\par\nobreak

（1）某盒中有4个红球、3个白球，依次不放回抽两个。已知第二个是白球，求第一个是红球的概率。

（2）在本节生产线模型中，已知抽到的产品合格，求它来自乙线的概率。

（3）盒甲有2个红球、1个白球，盒乙有1个红球、3个白球。等可能地选一盒，再从中随机取一球。已知取出红球，求所选为盒甲的概率。若改为从全部7个球中等可能抽一个，条件概率如何改变？

（4）设\(P(B)=0.3\)、\(P(A\mid B)=0.8\)、\(P(A\mid B^c)=0.2\)。求\(P(A)\)、\(P(B\mid A)\)及\(P(B\mid A^c)\)。

（5）（据2024年天津卷改编）从包含“整地做畦”和“田间灌溉”的5个劳动技能项目中等可能选3项。求选中前一项的概率，以及已选中前一项时也选中后一项的概率。
\end{exercise}

\FloatBarrier
\newpage
\subsection{信息与分割*}

将\(\Omega\)分成两两不交且非空的\(B_1,\ldots,B_m\)，就得到一个{\heiti 分割}。若只记录样本点属于哪个\(B_i\)，同一组内的结果在当前信息下便无法区分。对\(\Omega\)的每个样本点都单独记录，给出最细的分割；只记录必然事件，则没有获得新的区分。

例如，抛两次硬币时，“正面总次数”给出分割
\[
 \{HH\},\quad\{HT,TH\},\quad\{TT\};
\]
“第一次的面”给出\(\{HH,HT\},\{TH,TT\}\)。这两种记录各保留一部分信息，前者能确定总次数，后者能确定先后位置。若同时记录两者，就能区分四个结果。

若分割\(\mathcal Q\)中的每一组都包含于分割\(\mathcal P\)的某一组，称\(\mathcal Q\)比\(\mathcal P\){\heiti 更细}。从粗记录获得更细记录，原来的条件概率可以继续按组分解。例如\(B=C_1\cup\cdots\cup C_r\)是互斥并，则
\[
 P(A\mid B)=\sum_j P(C_j\mid B)P(A\mid C_j).
\]
这是在条件模型\(B\)内再次使用全概率公式。后面按信息分组求条件平均，也将遵循相同的结构。

\begin{exercise}
掷两枚骰子，设36个有序点数对等可能。分别考虑记录“点数和”“第一枚点数”“两个点数的有序对”所形成的分割。比较其粗细关系，并求已知点数和为8时第一枚大于3的概率。
\end{exercise}

\FloatBarrier
\newpage
\section{独立性}
\label{sec:prob-independence}

条件概率研究信息怎样改变概率。若知道一个事件发生，另一个事件的概率保持不变，就有了独立性的直观含义。下面用不需要相除的等式来定义它。

\subsection{事件的独立性}

\begin{definition}
若事件\(A,B\)满足
\begin{equation}\label{eq:prob-event-independence}
 P(A\cap B)=P(A)P(B),
\end{equation}
则称\(A,B\){\heiti 相互独立}。
\end{definition}

当\(P(B)>0\)时，此条件等价于\(P(A\mid B)=P(A)\)。在独立抛掷的两枚公平骰子中，第一枚出现偶数与第二枚大于4独立；若两事件都只涉及第一枚，则需要按实际事件关系计算，不能仅凭叙述中有两个事件就使用乘法。

独立性还保留于补事件。由
\[
 P(A\cap B^c)=P(A)-P(A\cap B)=P(A)(1-P(B)),
\]
得到\(A,B^c\)独立；同理\(A^c,B\)及\(A^c,B^c\)也独立。

互斥描述不能同时发生，独立描述概率的乘积分解。两个正概率事件若互斥，交事件概率为零，而概率乘积为正，因此不独立。

\subsection{多个事件的独立性}

\begin{definition}
若\(A_1,\ldots,A_n\)中任意非空子组均满足
\[
 P(A_{i_1}\cap\cdots\cap A_{i_r})=
 P(A_{i_1})\cdots P(A_{i_r}),
\]
则称这些事件{\heiti 相互独立}。若仅每两个事件独立，称为两两独立。
\end{definition}

两两独立不足以保证相互独立。公平抛两次硬币，令\(A\)为第一次正面，\(B\)为第二次正面，\(C\)为两次相同。三者的概率均为\(1/2\)，各两事件交集的概率均为\(1/4\)，但
\[
 P(A\cap B\cap C)=\frac14\ne\frac18.
\]
知道\(A,B\)同时发生，已经确定了\(C\)发生；各自单独考察时却看不出这种三者关系。

相互独立也允许将其中任意事件换成补事件。例如，若\(A,B,C\)相互独立，则
\[
 P(A\cap B\cap C^c)=P(A\cap B)-P(A\cap B\cap C)
 =P(A)P(B)P(C^c).
\]
对其他补事件重复使用这种分解，就得到一般情形。

\subsection{重复试验}

一次试验只记录成功与失败，成功概率为\(p\)，称为一次{\heiti Bernoulli试验}。独立重复\(n\)次且各次的\(p\)相同，就得到独立同机制的重复试验。

指定一条含\(k\)次成功、\(n-k\)次失败的结果序列，其概率为\(p^k(1-p)^{n-k}\)。所有恰有\(k\)次成功的序列共有\(\binom nk\)条，因而
\begin{equation}\label{eq:prob-bernoulli-count}
 P(\text{恰成功}k\text{次})=\binom nkp^k(1-p)^{n-k}.
\end{equation}
至少成功一次的概率为\(1-(1-p)^n\)，直接由“全部失败”的补事件得到。

独立性是试验机制的一项模型条件。来自同一有限盒子的不放回抽取，会改变后续盒中构成；独立重复测量则需要按测量机制判断，不能由次数相同代替判断。

\begin{example}[据2022年全国乙卷理科改编]
某棋手与甲、乙、丙各比赛一盘，各盘结果相互独立，获胜概率分别为\(0<p_1<p_2<p_3<1\)。记三盘中至少出现一段连续两盘获胜的概率为\(q\)。比较三种不同的第二盘对手，确定怎样安排使\(q\)最大。
\end{example}

\begin{exercise}
\mbox{}\par\nobreak

（1）独立重复试验5次，每次成功概率为\(0.4\)。求全部成功、至少成功一次及恰成功两次的概率。

（2）两个独立元件正常工作的概率分别为\(0.9,0.8\)。串联系统要求二者都正常，并联系统要求至少一个正常。求两种系统正常的概率。

（3）已知\(P(A)=P(B)=1/2\)，\(A,B\)独立。求\(P(A\cup B)\)及\(P(A\mid A\cup B)\)。观察获得新信息后，所用模型发生了什么变化。

（4）独立公平抛三次硬币，令\(A\)表示前两次相同，\(B\)表示后两次相同，\(C\)表示第一、三次相同。判断三事件是否两两独立、相互独立。

（5）（据2021年天津卷改编）甲、乙每次各猜一个谜语，猜对概率分别为\(5/6,3/5\)。每次两人的结果互不影响，各次活动也互不影响。一人猜对、另一人猜错时，猜对者获胜，否则平局。求甲一次获胜及三次活动中至少获胜两次的概率。
\end{exercise}

\FloatBarrier
\newpage
\subsection{条件信息下的独立性*}

在\(P(C)>0\)的条件模型中，若
\[
 P(A\cap B\mid C)=P(A\mid C)P(B\mid C),
\]
称\(A,B\)在条件\(C\)下独立。它描述的是当前条件模型中的乘积分解。

例如先等可能选择机器甲或乙，再用所选机器独立制造两件产品。甲的不合格率为\(0.1\)，乙为\(0.9\)。已知机器来源，两件是否不合格相互独立；不记录来源时，每件不合格的概率均为\(0.5\)，但两件都不合格的概率为
\[
 \frac12(0.1)^2+\frac12(0.9)^2=0.41\ne0.25.
\]
第一件不合格会提供关于机器来源的信息，也就改变了第二件不合格的概率。

反过来，独立事件在获得新信息后也可能不再独立。公平抛两次硬币，两个正面事件原先独立；已知至少一个正面时，每个正面事件的条件概率为\(2/3\)，同时发生的条件概率却为\(1/3\)。判断独立性需要明确在哪个概率模型中判断。

\begin{exercise}
在上述两台机器的模型中，已知第一件不合格，求第二件也不合格的概率；再求已知第一件合格时第二件不合格的概率。比较两者与事先的\(0.5\)。
\end{exercise}

\FloatBarrier
\newpage
\section{随机变量}
\label{sec:prob-random-variables}

同一次试验可以记录不同的数值。抛两次硬币，可以记录正面总数，也可以记录第一次是否正面。数值由实际结果确定，于是可以使用函数的语言来研究它。

\subsection{随机变量的定义}

\begin{definition}
在有限样本空间\(\Omega\)上，一个从样本点到实数的函数\(X:\Omega\to\mathbb R\)称为{\heiti 随机变量}。
\end{definition}

随机的是试验结果，函数\(X\)则把每个结果对应到一个确定的数。例如在两次抛币中定义正面次数\(Y\)，便有
\[
 Y(HH)=2,\quad Y(HT)=Y(TH)=1,\quad Y(TT)=0.
\]
\(\{Y=1\}\)表示事件\(\{HT,TH\}\)，\(\{Y\ge1\}\)表示\(\{HH,HT,TH\}\)。一般地，\(\{X\in A\}\)是所有满足\(X(\omega)\in A\)的样本点组成的事件，它的概率由原模型确定。

若对数值再作运算，\(g(X)\)仍是样本点的函数，也就是新的随机变量。两个变量的和\(X+Y\)、乘积\(XY\)，同样按每个样本点上的函数值定义。

\subsection{分布列}

\begin{definition}
有限随机变量\(X\)的不同可能取值为\(x_1,\ldots,x_m\)，记
\[
 p_j=P(X=x_j),\qquad j=1,\ldots,m.
\]
取值与相应概率组成的表称为\(X\)的{\heiti 分布列}。
\end{definition}

事件\(\{X=x_j\}\)分割样本空间，故\(p_j\ge0\)、\(\sum_jp_j=1\)。两次独立公平抛币的正面次数满足
\[
\begin{array}{c|ccc}
Y&0&1&2\\\hline
P&1/4&1/2&1/4
\end{array}
\]
分布列保留了各数值的概率，却不再逐一记录这些数值来自哪些样本点。不同试验中的变量可以有相同分布；同一空间中的两个变量也可以有相同分布而取值并不处处相同。

求\(g(X)\)的分布时，把产生同一新数值的概率合并。例如\(Z=(Y-1)^2\)只取\(0,1\)，两者的概率都为\(1/2\)。

\begin{example}
公平掷两枚骰子，\(X\)为较大点数。求\(X\)与\(6-X\)的分布列。
\end{example}

\newpage
\subsection{分布函数}

\begin{definition}
随机变量\(X\)的{\heiti 分布函数}定义为
\begin{equation}\label{eq:prob-cdf}
 F_X(x)=P(X\le x),\qquad x\in\mathbb R.
\end{equation}
\end{definition}

它记录从左侧累积到\(x\)的概率，因而单调不减，取值在\([0,1]\)中。对有限变量，在最小可能值左侧为0，在最大可能值及其右侧为1。例如上述\(Y\)满足
\[
 F_Y(x)=\begin{cases}
 0,&x<0,\\1/4,&0\le x<1,\\3/4,&1\le x<2,\\1,&x\ge2.
 \end{cases}
\]

\begin{figure}[!ht]
\centering
\includegraphics[width=.55\linewidth]{probability-cdf}
\caption{两次独立公平抛币的正面次数的分布函数。实心端点表示该处实际函数值。}
\end{figure}

在相邻取值之间，事件\(\{X\le x\}\)不变，图像便保持水平。由于定义包含端点，有限变量的分布函数在跳跃处取跳跃后的高度，并且右连续；图中的实心与空心端点正反映了这一点。

由事件分解，对\(a<b\)总有
\begin{equation}\label{eq:prob-cdf-interval}
 P(a<X\le b)=F_X(b)-F_X(a).
\end{equation}
对离散变量，跳跃大小给出该点的概率：\(P(X=a)=F_X(a)-F_X(a^-)\)，其中\(F_X(a^-)\)为左极限。有限模型中的这一累计表达，后面还可以用于连续取值的模型。

\begin{example}
利用上述\(F_Y\)，求\(P(Y<1)\)、\(P(Y\le1)\)与\(P(1\le Y\le2)\)，再从各跳跃高度恢复\(Y\)的分布列。
\end{example}

\subsection{联合分布}

研究两个变量的共同变化，需要知道它们同时取值的概率。若\(X\)取\(x_i\)、\(Y\)取\(y_j\)，记
\[
 p_{ij}=P(X=x_i,Y=y_j).
\]
这些概率组成{\heiti 联合分布列}，各项非负，总和为1。对同一行或列求和，得到各变量自身的分布，称为{\heiti 边际分布}。

在两次独立公平抛币中，令\(X\)记录第一次是否正面，\(Y\)记录正面总次数，便有
\[
\begin{array}{c|ccc|c}
 X\backslash Y&0&1&2&P(X=x)\\\hline
 0&1/4&1/4&0&1/2\\
 1&0&1/4&1/4&1/2\\\hline
 P(Y=y)&1/4&1/2&1/4&1
\end{array}
\]
已知\(Y=1\)，将这一列的概率除以列和\(1/2\)，就得到\(X\)在该条件下的分布。已知\(Y=2\)时，\(X=1\)的条件概率则为1。

由联合表也可以求\(X+Y\)等变量的分布：先标记每个格子的函数值，再把相同值对应的概率相加。这仍是样本结果合并的同一过程。

\subsection{随机变量的独立性}

\begin{definition}
若对全部可能取值均有
\begin{equation}\label{eq:prob-variable-independence}
 P(X=x_i,Y=y_j)=P(X=x_i)P(Y=y_j),
\end{equation}
称有限随机变量\(X,Y\)相互独立。
\end{definition}

对任意取值集合\(A,B\)，把上式在相应格子中求和，得到
\[
 P(X\in A,Y\in B)=P(X\in A)P(Y\in B).
\]
所以这个定义与事件独立性一致，也说明分别作用于\(X,Y\)的函数\(g(X),h(Y)\)仍然独立。

多个随机变量相互独立，要求任取有限子组及各自的取值，相应联合概率均等于各边际概率的乘积。连续取值时，则用各变量分别落入任意区间的事件满足同样的乘积分解来定义独立性。

两次独立抛币各自的正面指示量独立；第一次的正面指示量与总正面次数却不独立。即使给定两张相同的边际分布，也不能由此确定联合分布。例如\(X,Y\)均以\(1/2\)的概率取0和1，若\(Y=X\)，和只取0、2；若二者独立，和还以\(1/2\)的概率取1。

\begin{exercise}
\mbox{}\par\nobreak

（1）随机变量\(X\)在\(-1,0,1\)处的概率分别为\(1/4,1/2,1/4\)。写出\(X^2\)和\(2X+1\)的分布列及\(X\)的分布函数。

（2）从编号1至4的球中不放回抽两个，令\(X\)为第一次的数字，\(Y\)为第二次的数字。写出联合分布，求\(P(X+Y=5)\)，并判断独立性。

（3）在本节抛币联合表中，求\(X+Y\)的分布及\(P(X=1\mid Y\ge1)\)。

（4）设\(X,Y\)独立，各以\(1/2\)的概率取\(-1,1\)。求\(X+Y\)、\(XY\)的分布，判断\(X\)与\(XY\)是否独立。
\end{exercise}

\FloatBarrier
\newpage
\section{数学期望}
\label{sec:prob-expectation}

分布列列出了全部取值及其概率。有时我们先需要一个描述平均水平的数。各取值出现的可能性不同，平均时便应使用相应的概率作为权重。

\subsection{加权平均}

\begin{definition}
若有限随机变量\(X\)的分布列为\(P(X=x_i)=p_i\)，定义其{\heiti 数学期望}为
\begin{equation}\label{eq:prob-expectation}
 E(X)=\sum_i x_ip_i.
\end{equation}
\end{definition}

例如，一次试验的得分为\(0,2,5\)，概率依次为\(1/2,1/3,1/6\)，则期望为\(3/2\)。这个数可以不属于实际取值，它描述模型中的加权平均水平。

将同一取值对应的样本点展开，得到等价形式
\[
 E(X)=\sum_{\omega\in\Omega}X(\omega)P(\{\omega\}).
\]
因此求\(g(X)\)的期望，可以直接使用
\begin{equation}\label{eq:prob-expectation-function}
 E\bigl(g(X)\bigr)=\sum_i g(x_i)p_i,
\end{equation}
不必先把相同的\(g(x_i)\)合并成新的分布列。常值\(c\)的期望为\(c\)，而\(X\ge0\)时\(E(X)\ge0\)。有限取值的最大、最小值也夹住其期望。

\begin{example}
一次活动的奖金为0元、10元、50元，相应概率为\(0.7,0.25,0.05\)。求奖金期望。若每次参加需支付5元，求净收益期望，并说明期望与一次参加的实际收益分别表示什么。
\end{example}

\subsection{期望的线性性}

\begin{theorem}\label{thm:prob-linearity}
同一有限概率空间上的随机变量满足
\begin{equation}\label{eq:prob-expectation-linearity}
 E(aX+bY)=aE(X)+bE(Y).
\end{equation}
这里不要求\(X,Y\)独立。
\end{theorem}

\begin{proof}
在各样本点上展开：
\[
 \sum_\omega(aX(\omega)+bY(\omega))P(\{\omega\})
 =a\sum_\omega X(\omega)P(\{\omega\})
 +b\sum_\omega Y(\omega)P(\{\omega\}).\qedhere
\]
\end{proof}

连续使用此式即可求有限多个变量之和的期望。它允许先把总量拆成容易处理的部分，而不求总量的完整分布。例如两次掷骰的点数和的期望，为两枚点数期望之和；不放回抽出的各球数字之和，也可以逐次求期望后相加。

乘积的期望需要额外依据。若\(X,Y\)独立，由联合分布的乘积分解可得
\begin{equation}\label{eq:prob-independent-product}
 E(XY)=\sum_{i,j}x_iy_jp_iq_j=E(X)E(Y).
\end{equation}
一般变量不具有这一乘积性质，例如\(X\)以相同概率取\(-1,1\)时，\(E(X^2)=1\)，而\((EX)^2=0\)。

\subsection{指示变量}

对事件\(A\)，定义{\heiti 指示变量}
\[
 1_A(\omega)=\begin{cases}1,&\omega\in A,\\0,&\omega\notin A.\end{cases}
\]
由定义立即得到\(E(1_A)=P(A)\)。若\(N\)统计\(n\)个对象中满足条件的个数，将第\(j\)个对象是否满足条件记为\(I_j\)，就有
\begin{equation}\label{eq:prob-indicators}
 N=\sum_{j=1}^nI_j,\qquad E(N)=\sum_{j=1}^nP(I_j=1).
\end{equation}
这些指示变量可以相关，线性性仍然成立。

当\(X\)只取\(0,1,\ldots,m\)时，还可以按“至少达到第几层”拆分：
\[
 X=\sum_{r=1}^m1_{\{X\ge r\}},\qquad
 EX=\sum_{r=1}^mP(X\ge r).
\]
取值为\(k\)时，前\(k\)个指示变量为1，其余为0。这给出了用尾概率计算期望的方法，后面处理有次数上限的等待问题时会用到它。

例如将\(n\)封不同信件随机放入\(n\)个已写好收信人的信封，令\(N\)为放对的信件数。每封放对的概率为\(1/n\)，所以\(E(N)=1\)。若直接求\(N\)的分布，会涉及各种错放方式；期望只需要各封信的单个概率。

\begin{example}
从\(N\)个产品中不放回随机抽取\(n\)个，其中有\(M\)个不合格。分别按每次抽取、每个不合格产品是否入选定义指示变量，用两种分解求不合格品数的期望。
\end{example}

\subsection{条件平均}

已知事件\(B\)发生，将原分布换成条件分布，就得到条件下的平均。对\(P(B)>0\)，定义
\[
 E(X\mid B)=\sum_i x_iP(X=x_i\mid B).
\]
设\(B_1,\ldots,B_m\)构成正概率事件的分割，按样本点先分组求和，得到{\heiti 全期望公式}
\begin{equation}\label{eq:prob-total-expectation}
 E(X)=\sum_jP(B_j)E(X\mid B_j).
\end{equation}
各组平均需要按各组在总体中的概率加权。

具体地，条件平均乘以组的概率后，便恢复这一组对总平均的贡献：
\[
 \sum_jP(B_j)E(X\mid B_j)
 =\sum_j\sum_i x_iP(\{X=x_i\}\cap B_j)
 =\sum_i x_iP(X=x_i)=EX.
\]

例如甲、乙两类订单比例为\(0.4,0.6\)，平均处理时间为2分钟、5分钟，全部订单的平均处理时间就是\(0.4\cdot2+0.6\cdot5=3.8\)分钟。这里分组平均及其权重共同决定总平均。

\begin{example}[据2024年新课标II卷改编]
甲、乙组成投篮队，每次命中概率分别为\(0<p<q<1\)，各次结果独立。一人先投3次，全部未中则全队得0分；至少命中一次，则另一人再投3次，每次命中得5分，全队只计第二阶段得分。分别求两种出场顺序下的期望得分，并确定哪种顺序的期望较大。
\end{example}

\begin{exercise}
\mbox{}\par\nobreak

（1）随机变量\(X\)以\(1/3,1/2,1/6\)的概率取\(-1,1,3\)。求\(E(X)\)、\(E(X^2)\)及\(E(2X-3)\)。

（2）把五张标号1至5的卡片随机排在相应编号的位置上，求位置正确的卡片数的期望。

（3）十人各等可能选择五个窗口之一办理业务，各人选择独立。用指示变量求无人选择的窗口数的期望。

（4）从甲、乙两校各随机抽一人，分别以\(0.8,0.75\)的概率答对某题，答题相互独立。求答对人数的分布与期望，再用线性性核对。

（5）设有\(n\)个相互独立的生日，均匀分布在365天中。求生日相同的人对数的期望；一对按无序的两人识别。
\end{exercise}

\FloatBarrier
\newpage
\subsection{期望作为线性泛函*}

在\(\Omega=\{\omega_1,\ldots,\omega_N\}\)上，全体实值函数按逐点加法与数乘组成线性空间\(V\)。用函数值组成的列向量，可以把它与\(\mathbb R^N\)对应。

期望是映射\(E:V\to\mathbb R\)。式\eqref{eq:prob-expectation-linearity}说明它是线性映射，从线性空间到数域的线性映射称为{\heiti 线性泛函}。记常值函数为\(\mathbf1\)，则\(E(\mathbf1)=1\)。每个变量有分解
\[
 X=E(X)\mathbf1+\bigl(X-E(X)\mathbf1\bigr),
\]
第二项的期望为零。因此
\[
 V=\operatorname{span}\{\mathbf1\}\oplus\ker E.
\]
它把平均水平与零均值变化分开。有限空间中，指示函数\(1_{\{\omega_i\}}\)组成一组基，且
\[
 X=\sum_iX(\omega_i)1_{\{\omega_i\}},\qquad
 E(X)=\sum_iX(\omega_i)E(1_{\{\omega_i\}}).
\]
基函数的期望就是各样本点的概率，期望的线性性质由此与最初的权重分配接在一起。

\begin{exercise}
\mbox{}\par\nobreak

（1）求\(\ker E\)的维数。在三个等可能样本点的空间上写出它的一组基。

（2）设线性映射\(L:V\to\mathbb R\)满足\(L(\mathbf1)=1\)，且非负函数的像均非负。令\(p_i=L(1_{\{\omega_i\}})\)，证明它们构成一组概率权重，且\(L\)就是相应的期望。
\end{exercise}

\FloatBarrier
\newpage
\section{方差}
\label{sec:prob-variance}

平均水平相同的变量，取值可以集中，也可以分散。因为\(E(X-EX)=0\)，正负偏差直接相加不能反映波动。平方偏差保留偏离的大小，并使各项非负。

\subsection{平方偏差}

\begin{definition}
随机变量\(X\)的{\heiti 方差}与{\heiti 标准差}分别定义为
\[
 \operatorname{Var}X=E\bigl((X-EX)^2\bigr),\qquad
 \sigma_X=\sqrt{\operatorname{Var}X}.
\]
\end{definition}

展开平方并使用线性性，得到
\begin{equation}\label{eq:prob-variance-computation}
 \operatorname{Var}X=E(X^2)-(EX)^2.
\end{equation}
方差为零，当且仅当所有正概率结果上的\(X\)都等于同一常数。平移不改变波动，伸缩把平方偏差乘以伸缩倍数的平方：
\begin{equation}\label{eq:prob-variance-affine}
 E(aX+b)=aEX+b,\qquad \operatorname{Var}(aX+b)=a^2\operatorname{Var}X.
\end{equation}

对任意常数\(c\)，写\(X-c=(X-EX)+(EX-c)\)，交叉项期望为零，故
\begin{equation}\label{eq:prob-square-error}
 E(X-c)^2=\operatorname{Var}X+(EX-c)^2.
\end{equation}
所以用一个常数近似\(X\)，使平均平方误差最小的常数恰为\(EX\)。平均与方差在这里共同描述了一个最优近似问题。

若\(X\)以某个长度单位记录，方差使用该单位的平方，标准差则与\(X\)使用同一单位。因此用标准差描述典型偏离的尺度，更便于与原来的取值比较。

\begin{example}
\(X\)以相同概率取\(-1,1\)，\(Y\)以相同概率取\(-3,3\)。比较它们的期望、方差与标准差，再比较\(2X+5\)和\(Y+5\)。
\end{example}

\begin{example}[据2019年浙江卷改编]
设\(0<a<1\)，\(X\)等可能取\(0,a,1\)。研究\(a\)增大时\(\operatorname{Var}X\)怎样变化，并求其最小值。
\end{example}

\subsection{协方差}

研究\(aX+bY\)的波动时，平方展开还会产生\((X-EX)(Y-EY)\)这一交叉项。两者偏差同向时乘积为正，反向时为负；它的平均反映两个变量的共同变化。

\begin{definition}
随机变量\(X,Y\)的{\heiti 协方差}定义为
\[
 \operatorname{Cov}(X,Y)=E\bigl((X-EX)(Y-EY)\bigr)
 =E(XY)-EX\,EY.
\]
\end{definition}

由平方展开得到
\begin{equation}\label{eq:prob-sum-variance}
 \operatorname{Var}(aX+bY)=a^2\operatorname{Var}X+b^2\operatorname{Var}Y
 +2ab\operatorname{Cov}(X,Y).
\end{equation}
协方差对各变量分别线性。推广到多个变量，可将和的方差写为各方差之和，加上所有两两协方差的两倍。

独立时，式\eqref{eq:prob-independent-product}给出协方差为零，因而独立变量的和满足方差相加。反过来不成立。令\(X\)等可能取\(-1,0,1\)，\(Y=X^2\)，则\(EX=E(XY)=0\)，协方差为零；但\(Y\)由\(X\)完全确定，且\(P(Y=0\mid X=0)=1\ne P(Y=0)\)。

\begin{example}
两种测量方法的误差为\(X,Y\)，均值都为0，方差分别为1、4，协方差为0。将结果按权重\(a,1-a\)平均，误差为\(aX+(1-a)Y\)。求使该误差方差最小的权重，并与等权平均比较。
\end{example}

\subsection{相关系数}

当两方差都为正时，定义{\heiti 相关系数}
\begin{equation}\label{eq:prob-correlation}
 \rho_{XY}=\frac{\operatorname{Cov}(X,Y)}{\sigma_X\sigma_Y}.
\end{equation}
它把变量的单位与尺度除去。若\(\rho=0\)，称两变量{\heiti 不相关}。

在有限模型中，对各样本点上的数
\(\sqrt{p_i}(X(\omega_i)-EX)\)、\(\sqrt{p_i}(Y(\omega_i)-EY)\)使用Cauchy不等式，得到
\[
 |\operatorname{Cov}(X,Y)|^2\le\operatorname{Var}X\operatorname{Var}Y,
\]
故\(-1\le\rho\le1\)。等号成立恰好要求两列中心化数值成比例，即在正概率结果上\(Y=aX+b\)。\(a>0\)时\(\rho=1\)，\(a<0\)时\(\rho=-1\)。它度量的是这种线性共同变化。

换用不同的计量单位不会改变相关程度。由协方差与标准差的伸缩关系，当\(a,c\ne0\)时，
\[
 \rho_{aX+b,\,cY+d}=\frac{ac}{|ac|}\rho_{XY}.
\]
正比例变换和平移保持相关系数；只把一个变量的方向反转，相关系数便改变符号。

\newpage
\subsection{尾概率估计}

分布未知时，有时仍可以由少量数字特征控制偏离的概率。

\begin{theorem}[Markov不等式]
对非负随机变量\(Y\)及\(a>0\)，有
\begin{equation}\label{eq:prob-markov-bound}
 P(Y\ge a)\le\frac{EY}{a}.
\end{equation}
\end{theorem}

\begin{proof}
逐点有\(Y\ge a1_{\{Y\ge a\}}\)，取期望即得。
\end{proof}

将它用于非负变量\((X-EX)^2\)，就得到
\begin{theorem}[Chebyshev不等式]\label{thm:prob-chebyshev}
若\(X\)具有有限方差，则对\(t>0\)，
\begin{equation}\label{eq:prob-chebyshev}
 P(|X-EX|\ge t)\le\frac{\operatorname{Var}X}{t^2}.
\end{equation}
\end{theorem}

上述证明先在有限模型中建立；对后面具有相应有限期望或方差的离散、连续模型，同样的非负比较与期望运算仍成立。例如均值为100、标准差为2的变量，偏离均值至少10的概率不超过\(4/100=0.04\)。这个估计只利用方差，知道更多分布信息时可以进一步精确计算。

\begin{example}
\(X\)在\(-2,0,2\)处的概率分别为\(1/8,3/4,1/8\)。分别求\(P(|X|\ge1)\)、\(P(|X|\ge2)\)的实际值与Chebyshev上界。比较两种情形，判断何时界能取到等号。
\end{example}

\begin{exercise}
\mbox{}\par\nobreak

（1）\(X\)以\(1/4,1/2,1/4\)的概率取0、1、2。求\(EX\)、\(\operatorname{Var}X\)，并求使\(E(X-c)^2\)最小的\(c\)。

（2）设\(X,Y\)独立，期望分别为2、3，方差分别为1、4。求\(E(2X-Y)\)、\(\operatorname{Var}(2X-Y)\)及\(E(XY)\)。

（3）设\(\operatorname{Var}X=4\)、\(\operatorname{Var}Y=9\)、\(\operatorname{Cov}(X,Y)=-3\)。求相关系数及\(\operatorname{Var}(X+Y)\)。

（4）某非负变量的期望为3，求\(P(X\ge12)\)的上界。若还知道方差为1，使用Chebyshev不等式重新估计并比较。

（5）构造两个均值均为0、方差均为1的变量，使其和的方差分别为0、2、4。说明各构造中的相关关系。

（6）\(n\)次独立测量的误差均值为0、方差为9。求误差总和与平均误差的方差，并用Chebyshev界确定使平均误差绝对值至少1的概率不超过\(0.05\)的充分测量次数。
\end{exercise}

\FloatBarrier
\newpage
\subsection{平方误差的几何解释*}

保留有限样本空间中正概率的样本点，令其权重为\(p_i>0\)。在全体实值函数组成的空间上定义
\begin{equation}\label{eq:prob-weighted-inner-product}
 \langle X,Y\rangle=E(XY)=\sum_i p_iX(\omega_i)Y(\omega_i).
\end{equation}
线性、对称性由求和直接得到；正权重又保证\(\langle X,X\rangle\ge0\)，且等号只在\(X=0\)时成立，所以这是一个内积。这里度量所用的权重由概率模型确定。

常值函数\(\mathbf1\)的长度为1，且\(\langle X,\mathbf1\rangle=EX\)。因此\(EX\mathbf1\)正是\(X\)在常值方向上的正交投影，中心化残差满足
\[
 \langle X-EX\mathbf1,\mathbf1\rangle=0,\qquad
 \|X-EX\mathbf1\|^2=\operatorname{Var}X.
\]
式\eqref{eq:prob-square-error}成为勾股关系。协方差是中心化变量的内积，相关系数则是它们夹角的余弦。前面的平方展开与Cauchy不等式，由此具有共同的几何解释。

进一步，若只能用\(Y\)的仿射函数\(aY+b\)近似\(X\)，就要把\(X\)投影到\(\operatorname{span}\{\mathbf1,Y\}\)。设\(\operatorname{Var}Y>0\)，先匹配均值，得到\(b=EX-aEY\)；剩下的误差为
\[
 E\bigl[(X-EX)-a(Y-EY)\bigr]^2
 =\operatorname{Var}X-2a\operatorname{Cov}(X,Y)+a^2\operatorname{Var}Y.
\]
对\(a\)配方，最优近似为
\[
 \widehat X=EX+\frac{\operatorname{Cov}(X,Y)}{\operatorname{Var}Y}(Y-EY).
\]
当\(\operatorname{Var}X>0\)时，其最小均方误差为
\[
 E(X-\widehat X)^2=\operatorname{Var}X(1-\rho_{XY}^2).
\]
相关系数的平方由此刻画线性近似所解释的方差比例。后面的回归将对观测数据建立相应的投影。

\subsection{条件期望的投影解释*}

令\(\mathcal P=\{B_1,\ldots,B_m\}\)为正概率事件的分割。只能记录所属组别时，可使用的变量在每组内为常数，它们组成子空间
\[
 W=\left\{\sum_j c_j1_{B_j}:c_j\in\mathbb R\right\}.
\]
定义随机变量\(M\)，使它在\(B_j\)内的值为\(m_j=E(X\mid B_j)\)，记为\(M=E(X\mid\mathcal P)\)。这是把各组的条件平均组成一个新的变量。

令\(R=X-M\)。每组内残差的条件平均为零，所以对\(H=\sum_jh_j1_{B_j}\in W\)，
\[
 E(RH)=\sum_jh_jP(B_j)E(R\mid B_j)=0.
\]
因此\(X=M+R\)给出相对于\(W\)的正交分解，并对全部\(H\in W\)满足
\begin{equation}\label{eq:prob-conditional-projection}
 E(X-H)^2=E(X-M)^2+E(M-H)^2.
\end{equation}
在只能使用组别信息的范围内，条件平均使平方误差最小。

\begin{example}
四个样本点的概率依次为\(1/10,2/10,3/10,4/10\)，\(X\)的值依次为1、3、4、8。只记录样本点属于前两个点还是后两个点，求条件平均组成的变量\(M\)，并验证残差与两个组的指示变量都正交。
\end{example}

全期望公式给出\(EM=EX\)。再把\(X-EX\)分成\(X-M\)与\(M-EX\)，得到
\begin{equation}\label{eq:prob-total-variance}
 \operatorname{Var}X=\sum_jP(B_j)\operatorname{Var}(X\mid B_j)
 +\operatorname{Var}M.
\end{equation}
总波动分成组内波动的平均与组间平均的波动。后面的分组数据汇总及回归仍将使用这一平方误差结构。

若分割\(\mathcal Q\)比\(\mathcal P\)更细，记相应条件平均为\(M_{\mathcal Q},M_{\mathcal P}\)。粗分组下的变量也属于细分组的子空间，将\(H=M_{\mathcal P}\)代入细分组的正交分解，得到
\[
 E(X-M_{\mathcal P})^2
 =E(X-M_{\mathcal Q})^2+E(M_{\mathcal Q}-M_{\mathcal P})^2.
\]
新增信息所减少的最小平方误差，恰是两种条件平均之差的平方长度。

\begin{exercise}
\mbox{}\par\nobreak

（1）四个等可能样本点上，\(X\)的值依次为1、3、4、8。按前两个点与后两个点分组，求\(E(X\mid\mathcal P)\)，并核对总方差分解。

（2）在（1）的模型中，把后两个点进一步分别记录，求新增信息使最小平方误差减少多少，并与两种条件平均之差的平方长度比较。

（3）若分割\(\mathcal Q\)比\(\mathcal P\)更细，用正交分解证明\(E(E(X\mid\mathcal Q)\mid\mathcal P)=E(X\mid\mathcal P)\)。
\end{exercise}

\FloatBarrier
\newpage
\section{离散分布}
\label{sec:prob-discrete}

前面已建立分布、期望和方差的共同语言。现在根据试验机制，研究几种反复出现的分布。模型的选择先由抽取方式、停止规则和所记录的变量决定。

\subsection{二项分布}

独立重复\(n\)次Bernoulli试验，每次成功概率为\(p\)，令\(X\)为成功次数。由式\eqref{eq:prob-bernoulli-count}，
\begin{equation}\label{eq:prob-binomial-distribution}
 P(X=k)=\binom nkp^k(1-p)^{n-k},\qquad k=0,\ldots,n.
\end{equation}
称\(X\)服从参数为\(n,p\)的{\heiti 二项分布}，记为\(X\sim B(n,p)\)。二项式定理给出全部概率之和为\((p+1-p)^n=1\)。

令第\(i\)次成功的指示变量为\(I_i\)，则\(X=\sum_iI_i\)。每个\(I_i\)的期望为\(p\)，又因\(I_i^2=I_i\)，方差为\(p-p^2\)。独立性使协方差为零，从而
\begin{equation}\label{eq:prob-binomial-characteristics}
 EX=np,\qquad \operatorname{Var}X=np(1-p).
\end{equation}
相同成功概率与各次独立是此模型的共同条件。若每次成功概率不同，期望仍可相加，但成功总数通常不再服从同一个二项分布。

当\(0<p<1\)时，相邻两项的比值为
\[
 \frac{P(X=k+1)}{P(X=k)}=\frac{n-k}{k+1}\frac{p}{1-p},
 \qquad 0\le k<n.
\]
比值大于1等价于\(k+1<(n+1)p\)，因此概率先增后减，可以由\((n+1)p\)的位置确定概率最大的取值。

\begin{example}[2018年全国III卷理科，题面整理]
某群体每位成员使用移动支付的概率均为\(p\)，各人选择独立。10名成员中使用移动支付的人数为\(X\)，已知\(\operatorname{Var}X=2.4\)、\(P(X=4)<P(X=6)\)，求\(p\)。
\end{example}

\subsection{超几何分布}

\(N\)个不同对象中有\(M\)个属于目标类，等可能地选出\(n\)个，令\(X\)为其中目标类个数。选出\(k\)个目标类时，还需从另\(N-M\)个中选出\(n-k\)个，所以
\begin{equation}\label{eq:prob-hypergeometric}
 P(X=k)=\frac{\binom Mk\binom{N-M}{n-k}}{\binom Nn},
 \quad \max(0,n-N+M)\le k\le\min(n,M).
\end{equation}
这称为{\heiti 超几何分布}。Vandermonde恒等式核对了分布列的归一化。

把抽取随机排成先后次序，各次抽到目标类的概率均为\(p=M/N\)，但相互不独立。使用指示变量仍有\(EX=np\)。对两个不同抽取位置，
\[
 E(I_iI_j)=\frac{M(M-1)}{N(N-1)},\qquad
 \operatorname{Cov}(I_i,I_j)=-\frac{p(1-p)}{N-1}.
\]
因而当\(N>1\)时，
\begin{equation}\label{eq:prob-hypergeometric-variance}
 \operatorname{Var}X=np(1-p)\frac{N-n}{N-1}.
\end{equation}
最后的因子体现了不放回抽样对波动的影响。抽取全部总体时目标类个数已确定；抽取比例很小时，这一因子接近1。

两个模型还有分布上的联系。固定抽取次数\(n\)，令\(N\to\infty\)、\(M/N\to p\)。在任何固定的抽取记录中，下一次抽到目标类的条件概率为“剩余目标数除以剩余总数”，它趋于\(p\)。含\(k\)个目标类的每条指定记录的概率便趋于\(p^k(1-p)^{n-k}\)，所以
\[
 P(X=k)\longrightarrow\binom nkp^k(1-p)^{n-k}.
\]
总体很大而抽样次数固定时，不放回模型因此趋近相同目标比例的二项模型。

\begin{example}
一批20件产品中有4件不合格，随机不放回抽5件。求不合格品数的分布、至少抽到一件不合格品的概率，以及均值与方差。若改为放回抽取，比较两个模型。
\end{example}

\subsection{几何分布}

独立重复成功概率为\(0<p<1\)的Bernoulli试验，直到第一次成功，令\(T\)为所需次数。事件\(\{T=k\}\)要求先失败\(k-1\)次再成功，所以
\begin{equation}\label{eq:prob-geometric}
 P(T=k)=(1-p)^{k-1}p,\qquad k=1,2,\ldots.
\end{equation}
称为{\heiti 几何分布}。这里模型具有可数无穷个可能取值，其概率和按数项级数理解。等比级数给出\(\sum_{k\ge1}P(T=k)=1\)，因此它确实形成概率分配。

记\(q=1-p\)。等待超过\(m\)次就是前\(m\)次均失败，故\(P(T>m)=q^m\)，并有
\[
 P(T>m+r\mid T>m)=\frac{q^{m+r}}{q^m}=q^r=P(T>r).
\]
这称为{\heiti 无记忆性}，它来自每次试验的独立和机制不变。

对非负可数离散变量，期望是取值乘概率的级数；级数收敛时便有有限期望。

几何分布的有限和为
\[
 \sum_{k=1}^n k p q^{k-1}=\frac{1-q^n}{p}-nq^n.
\]
由\(q^n\to0\)、\(nq^n\to0\)，可得
\begin{equation}\label{eq:prob-geometric-mean}
 E(T)=\frac1p.
\end{equation}
这个平均等待次数由无限取值的级数得到，不能把\(T\)的随机性直接替换成一个固定次数。

实际尝试常有次数上限。若第一次成功便停止，至多尝试\(m\)次，则实际次数为\(T_m=\min(T,m)\)。到达第\(m\)次只要求前\(m-1\)次失败，故\(P(T_m=m)=q^{m-1}\)。由前面用尾概率计算期望的公式，
\[
 E T_m=\sum_{r=1}^mP(T_m\ge r)
 =\sum_{r=1}^m q^{r-1}=\frac{1-q^m}{p}.
\]
在次数上限内成功的概率为\(1-q^m\)。这一概率与实际尝试次数的期望分别回答“能否成功”和“平均尝试多少次”。

\begin{exercise}
\mbox{}\par\nobreak

（1）某学生独立回答6道同类题，每题答对概率为\(3/4\)。求至少答对5题的概率及答对题数的期望与方差。

（2）从10个球中不放回抽4个，其中3个为红球。求红球数的分布与数字特征，并核对全部概率之和。

（3）每次投篮独立，命中率为\(0.4\)。求第一次命中发生在第3次、超过第5次的概率及平均等待次数。

（4）某人已连续失败5次，试验条件与独立性保持不变。求再等待超过3次的条件概率，并与尚未开始时比较。

（5）用有限和或第一章的级数求几何分布的\(E(T^2)\)，证明\(\operatorname{Var}T=(1-p)/p^2\)。

（6）说明下列变量适用的模型及参数：从有限产品中不放回抽取固定件数时的不合格数；固定次数独立抽检的不合格数；独立抽检直到首次发现不合格品时的抽检次数。

（7）（据2020年浙江卷改编）盒中有1个红球、1个绿球、2个黄球，每次随机取一球，不放回，取出红球即停。求已取黄球数的分布与期望，再用指示变量核对期望。这里的停止次数是否服从几何分布？
\end{exercise}

\FloatBarrier
\newpage
\subsection{Poisson近似*}

把\(n\)个球独立等可能地投入\(m\)个盒子，某个固定盒子中的球数服从\(B(n,1/m)\)。当球数、盒数都很大，且\(n/m\to\lambda>0\)时，每盒平均球数保持在\(\lambda\)附近，计数的分布会趋向什么？

一般地，设\(X_n\sim B(n,p_n)\)，且\(np_n\to\lambda>0\)。固定非负整数\(k\)，有
\[
 P(X_n=k)=\frac{(np_n)^k}{k!}
 \prod_{j=0}^{k-1}\left(1-\frac jn\right)
 (1-p_n)^n(1-p_n)^{-k}.
\]
当\(k=0\)时空乘积为1。由\(p_n\to0\)及\(\ln(1-p_n)/p_n\to-1\)，
\[
 n\ln(1-p_n)\to-\lambda,\qquad
 P(X_n=k)\longrightarrow e^{-\lambda}\frac{\lambda^k}{k!}.
\]
由指数级数，右侧各项之和为1。以这些概率定义的分布称为参数为\(\lambda\)的{\heiti Poisson分布}，记为\(\operatorname{Pois}(\lambda)\)。

例如200次独立试验，单次事件概率为\(0.005\)时，\(\lambda=1\)。完全不发生的精确概率为\(0.995^{200}\)，Poisson近似为\(e^{-1}\)。

独立的Poisson计数还可以相加。若\(X\sim\operatorname{Pois}(\lambda)\)、\(Y\sim\operatorname{Pois}(\mu)\)且二者独立，则
\begin{align*}
 P(X+Y=k)&=e^{-(\lambda+\mu)}\sum_{j=0}^k
 \frac{\lambda^j\mu^{k-j}}{j!(k-j)!}\\
 &=e^{-(\lambda+\mu)}\frac{(\lambda+\mu)^k}{k!}.
\end{align*}
所以\(X+Y\sim\operatorname{Pois}(\lambda+\mu)\)，参数随独立计数相加。

还可以把一次计数分成两类。设\(N\sim\operatorname{Pois}(\lambda)\)，给定\(N\)后，各次独立以概率\(r\in(0,1)\)归入甲类，其余归入乙类，记两类个数为\(A,B\)。对非负整数\(i,j\)，
\begin{align*}
 P(A=i,B=j)
 &=e^{-\lambda}\frac{\lambda^{i+j}}{(i+j)!}\binom{i+j}i r^i(1-r)^j\\
 &=\frac{e^{-\lambda r}(\lambda r)^i}{i!}
  \frac{e^{-\lambda(1-r)}[\lambda(1-r)]^j}{j!}.
\end{align*}
右侧分解成两个Poisson分布的乘积，因此\(A,B\)独立，参数分别为\(\lambda r,\lambda(1-r)\)。

\begin{exercise}
\mbox{}\par\nobreak

（1）计算上述试验中恰发生一次、至少发生两次的二项精确表达与Poisson近似。

（2）设\(X\sim\operatorname{Pois}(\lambda)\)。由指数级数求\(EX\)、\(E[X(X-1)]\)，证明\(\operatorname{Var}X=\lambda\)。
\end{exercise}

\newpage
\subsection{概率生成函数*}

独立计数相加时，要把各对取值的概率相乘，再按相同总数相加。这恰好是多项式乘法计算系数的规则，因此可以用多项式记录整个分布。设\(X\)只取\(0,1,\ldots,m\)，定义
\begin{equation}\label{eq:prob-pgf}
 G_X(s)=E(s^X)=\sum_{k=0}^mP(X=k)s^k.
\end{equation}
称为\(X\)的{\heiti 概率生成函数}。系数直接恢复分布，\(G_X(1)=1\)。有限多项式求导给出
\[
 G_X'(1)=EX,\qquad G_X''(1)=E[X(X-1)],
\]
故\(\operatorname{Var}X=G_X''(1)+G_X'(1)-(G_X'(1))^2\)。

若\(X,Y\)独立，由乘积期望性质，
\[
 G_{X+Y}(s)=E(s^Xs^Y)=G_X(s)G_Y(s).
\]
Bernoulli变量的生成函数为\(1-p+ps\)。独立的\(n\)次试验中，若成功概率分别为\(p_1,\ldots,p_n\)，成功总数\(S\)便有
\[
 G_S(s)=\prod_{i=1}^n(1-p_i+p_is).
\]
各\(p_i\)都为\(p\)时，成为\((1-p+ps)^n\)，恢复二项分布。线性性与方差相加给出\(ES=\sum_i p_i\)、\(\operatorname{Var}S=\sum_i p_i(1-p_i)\)。

在\(s=-1\)处，偶次幂保持符号，奇次幂改变符号。与\(G_X(1)=1\)联立，便得
\[
 P(X\text{为偶数})=\frac{1+G_X(-1)}2,\qquad
 P(X\text{为奇数})=\frac{1-G_X(-1)}2.
\]
\begin{exercise}
\mbox{}\par\nobreak

（1）独立抛掷两枚公平骰子，点数分别记为\(X,Y\)。写出各自生成函数，以乘积的系数求点数和为7的概率，再求点数和的期望与方差。若两次取值完全相同，为什么不能使用同一个乘积？

（2）三个事件相互独立，发生概率依次为\(1/5,1/3,2/5\)。用生成函数求恰有两个发生及发生个数为偶数的概率。

（3）\(X\)为有限非负整数变量。给定\(X=k\)时，把\(k\)个对象各自独立地以概率\(r\)保留，记留下的个数为\(Y\)。证明\(G_Y(s)=G_X(1-r+rs)\)，并求\(EY\)、\(\operatorname{Var}Y\)。
\end{exercise}

\FloatBarrier
\newpage
\section{连续分布}
\label{sec:prob-continuous}

时间、长度和测量误差常用连续实数记录。现在把模型从有限或可数取值扩展到连续取值；随机变量仍然是结果的数值记录，其概率需要由区间上的分配来描述。分布函数提供了前后共用的表达。

\subsection{累计概率}

对连续取值的变量，仍记\(F(x)=P(X\le x)\)。它继续表示累积到\(x\)为止的概率，区间概率仍由\(F(b)-F(a)\)给出。

例如在区间\([a,b]\)上均匀选一点，概率与子区间长度成正比。在\(a\le x\le b\)时，分布函数为\(F(x)=(x-a)/(b-a)\)；在区间左侧为0、右侧为1。
当区间向一个点缩小时，概率趋于零。每个单点的概率均为零，但整个\([a,b]\)的概率为1。

分布函数连续时没有单点的跳跃概率。下面进一步研究可以用面积权重描述的连续模型。

\subsection{概率密度}

前面的均匀选点模型中，区间内一段的概率等于长度乘\(1/(b-a)\)，可画成矩形面积。若概率分配不均匀，就用变化的高度代替水平线。

\begin{definition}
若非负函数\(f\)的曲线下总面积为1，且对任意区间\((u,v]\)有
\[
 P(u<X\le v)=\text{曲线 }y=f(x)\text{ 在 }u\text{ 到 }v\text{ 下方的面积},
\]
则称\(f\)为\(X\)的{\heiti 概率密度}，称\(X\)为{\heiti 连续型随机变量}。
\end{definition}

已学习积分时，可把上式写为\(\int_u^v f(x)\,dx\)，分布函数为从负无穷累积到\(x\)的面积。未调用积分时，先使用可以直接求面积的图形及给定的模型概率。

\([a,b]\)上均匀分布的密度为
\[
 f(x)=\begin{cases}1/(b-a),&a\le x\le b,\\0,&\text{其他}.
 \end{cases}
\]
记为\(X\sim U(a,b)\)。密度的高度乘以区间长度，才给出相应概率。例如\(U(0,0.1)\)的密度高度为10，但落在\([0,0.02]\)中的概率为\(10\cdot0.02=0.2\)。

均匀分布满足\(EX=(a+b)/2\)、\(\operatorname{Var}X=(b-a)^2/12\)，选讲中将用积分推导。

\subsection{正态分布}

高尔顿钉板把一系列微小的随机偏移累积起来：小球从上方落下，每碰到一层钉子，就向左或向右偏移一步。在理想模型中，各层的偏移相互独立，两个方向的概率均为\(1/2\)。经过\(n\)层，向右的次数\(X_n\sim B(n,1/2)\)，落入第\(k\)个槽的概率为\(\binom nk2^{-n}\)。中间槽对应的路径多，两侧槽对应的路径少。

\begin{figure}[!ht]
\centering
\begin{minipage}[c]{.42\linewidth}
\centering\includegraphics[width=\linewidth]{probability-galton}
\end{minipage}\hfill
\begin{minipage}[c]{.54\linewidth}
\centering\includegraphics[width=\linewidth]{probability-normal-limit}
\end{minipage}
\caption{左为八层理想钉板，槽内柱高按二项概率绘制；右为不同层数的标准化概率直方图与同一条正态密度曲线。右下阴影为\(-1\)至1之间的曲线下面积。}
\label{fig:prob-galton-normal}
\end{figure}
\FloatBarrier

落球数量增加，使各槽的频率更接近相应的二项概率；增加钉板层数，则改变了二项分布本身。由\(EX_n=n/2\)、\(\operatorname{Var}X_n=n/4\)，把横坐标改为
\[
 Z_n=\frac{X_n-n/2}{\sqrt n/2},
\]
便把各层数的中心与尺度统一。图\ref{fig:prob-galton-normal}右侧每根柱的面积等于相应概率，随着层数增加，柱形逐渐贴近同一条光滑曲线。这提示我们用连续密度描述大量独立小偏移的累积；这种近似的理论依据将在中心极限定理中建立。

这条曲线关于0对称，中间高、两侧逐渐降低。下面先定义这一分布族，用两个参数分别表示中心位置与分散尺度。

\begin{definition}
密度为
\begin{equation}\label{eq:prob-normal-density}
 f(x)=\frac1{\sigma\sqrt{2\pi}}
 \exp\left(-\frac{(x-\mu)^2}{2\sigma^2}\right),\qquad \sigma>0,
\end{equation}
的分布称为{\heiti 正态分布}，记为\(N(\mu,\sigma^2)\)。
\end{definition}

\newpage
图\ref{fig:prob-galton-normal}右侧的光滑曲线对应\(\mu=0,\sigma=1\)，称为标准正态密度；右下阴影面积就是相应变量取值在\(-1\)与1之间的概率。

这里引用高斯积分的归一化结论：式\eqref{eq:prob-normal-density}下的总面积为1。该分布的期望为\(\mu\)，方差为\(\sigma^2\)。改变\(\mu\)平移中心，改变\(\sigma\)调整分散尺度。一个图像近似钟形，可以提示考虑正态模型，模型是否合适还需由问题机制和数据支持。

若\(X\sim N(\mu,\sigma^2)\)，则{\heiti 标准化变量}
\[
 Z=\frac{X-\mu}{\sigma}
\]
服从标准正态分布\(N(0,1)\)。记其密度为\(\varphi\)，分布函数为\(\Phi\)。标准化把\(X\)区间换成相应的\(Z\)区间，所以
\begin{equation}\label{eq:prob-normal-standardization}
 P(a<X\le b)=\Phi\left(\frac{b-\mu}{\sigma}\right)
 -\Phi\left(\frac{a-\mu}{\sigma}\right).
\end{equation}
标准化结论可由区间面积的横向伸缩得到，后面的密度换元也会给出证明。

对称性给出\(\Phi(-z)=1-\Phi(z)\)。以下数值供本章使用：
\[
\begin{array}{c|rrrrrr}
z&0&1&1.645&1.96&2&2.576\\\hline
\Phi(z)&0.5000&0.8413&0.9500&0.9750&0.9772&0.9950
\end{array}
\]
由此得到\(P(|Z|\le1)\approx0.6827\)、\(P(|Z|\le2)\approx0.9545\)、\(P(|Z|\le3)\approx0.9973\)。这些是模型概率，并非每组观测数据必须满足的固定比例。

\begin{example}
某加工尺寸用\(N(100,4)\)描述，单位为毫米。求尺寸在98至102毫米内、超过104毫米及低于98毫米的概率。若合格范围为96至104毫米，估计独立抽检100件中的合格件数期望。
\end{example}

\begin{exercise}
\mbox{}\par\nobreak

（1）某班车每10分钟到站一次。把乘客到站相位视为在一个周期内均匀，求等待时间超过6分钟的概率及等待时间的均值。

（2）\(X\sim N(50,9)\)。求\(P(47<X<53)\)、\(P(X>55.88)\)，并求以50为中心、包含约95\%概率的区间。

（3）某密度在\([0,1]\)上为\(f(x)=cx\)，其他地方为0。用三角形面积求\(c\)及\(P(X\le1/2)\)。

（4）两个正态模型具有相同均值，标准差分别为1、2。比较两者落在均值两侧各1单位区间内的概率，并说明标准化中的分母。
\end{exercise}

\FloatBarrier
\newpage
\subsection{连续型变量的期望*}

先考察有界区间上的连续密度与连续函数\(g\)。把区间细分，在每一小段用\(g(\xi_i)\)近似该段取值，并用\(f(\xi_i)\Delta x_i\)近似概率，加权和便趋于
\begin{equation}\label{eq:prob-continuous-expectation}
 E(g(X))=\int g(x)f(x)\,dx.
\end{equation}
这是离散加权平均的积分表达。在一般密度模型中，也用此式定义相应期望；对有正有负的函数，要求绝对值的积分有限，才取有限的期望。

若\(X\sim U(a,b)\)，则
\[
 EX=\frac1{b-a}\int_a^b x\,dx=\frac{a+b}2,
\]
并由\(E(X^2)=(a^2+ab+b^2)/3\)得到\(\operatorname{Var}X=(b-a)^2/12\)。

对无穷区间，积分按第五章的反常积分定义，分别考察两端极限。标准正态密度满足\(\varphi'(z)=-z\varphi(z)\)。因此
\[
 \int_0^\infty z\varphi(z)\,dz=\varphi(0)<\infty,
\]
再由对称性得到\(EZ=0\)。分部积分给出
\[
 \int_{-R}^R z^2\varphi(z)\,dz
 =[-z\varphi(z)]_{-R}^R+\int_{-R}^R\varphi(z)\,dz.
\]
因\(R\varphi(R)\to0\)且总面积为1，故\(E(Z^2)=1\)。若\(X=\mu+\sigma Z\)，则线性性与方差伸缩给出\(EX=\mu\)、\(\operatorname{Var}X=\sigma^2\)。

\subsection{变量变换*}

设\(X\)在区间\(I\)内取值，有密度\(f_X\)，\(Y=g(X)\)。先将事件\(\{Y\le y\}\)写回\(X\)的取值条件，再求分布函数，就能得到变换后的分布。

若\(g\)在\(I\)上严格递增且可微，导数不为零，记反函数为\(h\)，则
\[
 F_Y(y)=F_X(h(y)),\qquad f_Y(y)=f_X(h(y))h'(y).
\]
严格递减时，累积方向反转，得到统一公式
\begin{equation}\label{eq:prob-density-transform}
 f_Y(y)=f_X(h(y))|h'(y)|.
\end{equation}
它也说明标准化时为何需要相应的密度缩放，概率面积由横向伸缩与高度调整共同保留。

非单调变换则需要合并分支。例如\(X\sim U(-1,1)\)、\(Y=X^2\)，对\(0\le y\le1\)，
\[
 P(Y\le y)=P(-\sqrt y\le X\le\sqrt y)=\sqrt y.
\]
故\(f_Y(y)=1/(2\sqrt y)\)，\(0<y<1\)，其他地方为0。正、负两段\(X\)值都对同一个\(Y\)值作出贡献；将密度直接代入\(x=\sqrt y\)会遗漏区间伸缩与另一分支。

\begin{exercise}
\mbox{}\par\nobreak

（1）设\(X\sim U(-1,1)\)，求\(E(X^2)\)、\(E|X|\)，再用\(X^2\)的密度核对第一问。

（2）设\(Z\sim N(0,1)\)，求\(|Z|\)的分布函数、密度与期望。

（3）设\(X\)在\([0,1]\)上密度为\(2x\)。求\(Y=1-X\)的密度，并求\(EX\)、\(\operatorname{Var}X\)。
\end{exercise}

\FloatBarrier
\newpage
\section{数据的获取}
\label{sec:prob-sampling}

给定模型可以研究可能出现的数据。实际统计往往从另一端开始：面对尚不完全了解的总体，怎样获得有用的数据？抽样方式决定哪些对象可能入选，也决定随后概率计算与统计判断的依据。

\subsection{总体与样本}

\begin{definition}
研究所涉及的全部对象或其测量值组成{\heiti 总体}；实际选取并观察的一部分对象或其测量值组成{\heiti 样本}，观测个数称为{\heiti 样本量}。
\end{definition}

研究一批产品的不合格比例，总体是这批产品，记录变量为是否不合格。研究某种装置的测量误差，总体也可以由在固定条件下反复测量的概率分布描述。前者通常有一个明确的有限对象集合，后者主要由试验机制与分布模型确定。

总体比例、总体均值等描述总体的量称为{\heiti 参数}。样本中的比例、均值是由实际数据计算的量。取样以前，它们随样本变化；取样以后，它们成为具体的观测值。后面的抽样分布将进一步研究这种区别。

明确总体还包括明确时间、范围和测量方法。例如调查某校学生的每周锻炼时间，需要确定研究哪一学期、怎样计算锻炼时间，以及缺少记录时怎样处理。不同的记录规则可能形成不同的数据。

\subsection{简单随机抽样}

\begin{definition}
从\(N\)个不同对象中不放回抽取\(n\)个，若每个\(n\)元子集都具有概率\(1/\binom Nn\)，称为{\heiti 简单随机抽样}。
\end{definition}

可以将对象编号，用随机顺序取前\(n\)个，或均匀抽取不同编号。每个对象的入选概率为
\[
 \frac{\binom{N-1}{n-1}}{\binom Nn}=\frac nN.
\]
但仅有单个对象入选概率相同，还不足以保证所有子集等可能。例如四人中只在\(\{1,2\}\)、\(\{3,4\}\)两组之间等可能选择，每人的入选概率都是\(1/2\)，却没有给其他二人组任何机会。

有限总体中不放回抽样的各次结果相互关联。若采用有放回的独立抽取，或直接假定观测来自同一模型的独立重复试验，就得到另一种抽样模型；使用均值方差等公式时应按实际机制选取。

\subsection{分层随机抽样}

若总体由不同组别组成，可以先按有意义的特征分层，再在各层内随机抽取。设第\(h\)层有\(N_h\)个对象，样本量为\(n_h\)，总体量为\(N=\sum_hN_h\)。按比例分配时，令\(n_h/n\)尽量接近\(N_h/N\)。

若各层样本均值为\(\bar x_h\)，估计全体平均时应使用总体权重：
\begin{equation}\label{eq:prob-stratified-estimate}
 \widehat\mu=\sum_h\frac{N_h}{N}\bar x_h.
\end{equation}
比例分配时，它与全部样本直接计算的均值一致；非比例分配时，样本中的各层比例已经改变，需用上式恢复总体的构成。

例如两个年级分别有900人、100人，各抽20人，样本平均得分分别为60、90。按总体构成估计的平均为\(0.9\cdot60+0.1\cdot90=63\)，直接合并40人的平均则为75。后一个数描述当前样本的构成，前一个数用于恢复原总体的权重。

分层使各类对象有计划地进入样本，是否能降低估计的波动，还与层内差异和分配方式有关。本节先建立抽样机制与加权依据。

\subsection{研究设计}

随机抽样用于从目标总体取得样本；随机分组用于把参加者分配给不同处理。两者承担不同任务。随机抽样支持从样本向总体的推断，随机分组则帮助比较处理效果，减少原有差异对比较的系统影响。

观察性研究记录原有状态，例如比较自愿参加训练与未参加训练者的成绩。两组还可能在基础、时间投入等方面不同。对照试验则由研究者安排处理，在条件允许时随机分组，并尽量保持其他条件可比。即使观察到明显相关，因果解释仍需要研究设计提供依据。

自愿回应、只调查容易接触的人、部分对象未回应，以及测量口径不一致，都可能带来系统偏差。扩大同一种有偏样本，主要增加同一获取方式下的记录，不能自行恢复原来缺失的总体信息。

\begin{example}
某校准备调查全校学生的睡眠时间。比较以下方案：在图书馆晚间离馆者中调查；在全校学生名单中简单随机抽取；按年级分层随机抽取。说明各方案的目标总体、可能遗漏的对象，以及分层比例不同时怎样汇总平均值。
\end{example}

\begin{exercise}
\mbox{}\par\nobreak

（1）某总体分三层，人数为500、300、200，计划抽取100人。给出比例分配方案。若改为各层抽30人，三个样本均值为12、15、18，求总体均值的分层估计。

（2）为了研究某训练方法，从报名者中随机分为训练组与对照组。说明随机分组解决了什么问题；若希望结论适用于所有高中生，还需要考虑哪些取样信息？

（3）从编号1至20的对象中随机选一个起点\(1,\ldots,5\)，再每隔5个取一个。求每个对象的入选概率，并判断这是否是样本量4的简单随机抽样。

（4）（据2022年新高考II卷第19题改编）随机调查某地区100名某疾病患者，将年龄分入\([0,10),[10,20),\ldots,[80,90]\)，得到下表：
\[
\begin{array}{c|rrrrrrrrr}
 \text{组中值（岁）}&5&15&25&35&45&55&65&75&85\\\hline
 \text{频数}&1&2&12&17&23&20&17&6&2
\end{array}
\]
各组以组中值代表，估计患者的平均年龄及年龄在\([20,70)\)内的概率。若全地区患病率为\(0.1\%\)，年龄在\([40,50)\)内的人口占\(16\%\)，估计从这一年龄段随机选一人时其患病的概率，精确到\(0.0001\)。说明最后一问所用的总体与患者样本之间的关系。
\end{exercise}

\FloatBarrier
\newpage
\section{数据的描述}
\label{sec:prob-data-description}

先研究一组已经取得的数据。图表保留形态，数字特征概括位置与分散程度；选择怎样概括，需要联系所研究的变量和问题。本节使用教学数据
\begin{equation}\label{eq:prob-description-data}
 8,\ 9,\ 9,\ 10,\ 10,\ 11,\ 12,\ 15
\end{equation}
贯穿各种表示，单位暂不指定。

\subsection{分布图表}

取值类别较少时，可以用频数表记录每类个数，再除以总数得到频率。数值连续或类别很多时，可以按区间分组。若某组频率为\(f_i\)、组距为\(d_i\)，频率直方图的高度取\(f_i/d_i\)，使每个矩形的面积恰为该组频率。

例如将式\eqref{eq:prob-description-data}分入\([8,10)\)、\([10,12)\)、\([12,16)\)，频数依次为3、3、2。第三组虽有较大组距，其矩形面积仍为\(2/8\)。只比较柱高而忽略组距，会改变频率的解释。

\begin{figure}[!ht]
\centering
\includegraphics[width=.39\linewidth]{probability-histogram}\hfill
\includegraphics[width=.56\linewidth]{probability-histogram-exam}
\caption{左图标注各组频率；右图为零件直径数据，柱顶标注频率与组距之比。}
\label{fig:prob-histograms}
\end{figure}

分组会损失组内位置。需要精确计算时应保留原始数据；只有分组频数时，可以用组中值作近似，并将结果按近似量解释。两列成对数值则可用散点图保留对应关系，回归一节再详细研究。

\subsection{位置特征}

数据\(x_1,\ldots,x_n\)的{\heiti 样本均值}为\(\bar x=n^{-1}\sum_ix_i\)。将数据升序排列为\(x_{(1)}\le\cdots\le x_{(n)}\)，奇数个数据的中位数为中间一项，偶数个则取中间两项的平均；出现次数最多的值称为{\heiti 众数}，可以不唯一。

本书采用以下有限样本百分位数约定：对\(0<p<1\)，若\(np\)不是整数，第\(100p\)百分位数取\(x_{(\lceil np\rceil)}\)；若\(np=k\)为整数，则取\((x_{(k)}+x_{(k+1)})/2\)。第25、50、75百分位数分别记为\(Q_1,Q_2,Q_3\)，其中\(Q_2\)就是中位数。

本节数据的均值为10.5，中位数为10，众数为9和10；\(Q_1=9\)、\(Q_3=11.5\)。以\(Q_1,Q_3\)为箱体两端，在中位数处画线，再标记两端位置，就得到箱线图。本节采用两须延伸到最小值、最大值的简单画法。

\begin{figure}[!ht]
\centering
\includegraphics[width=.64\linewidth]{probability-boxplot}
\caption{同一组数据的箱线图，依次标出最小值、三个四分位数和最大值。}
\end{figure}

均值使\(\sum_i(x_i-c)^2\)最小，这由平方误差分解得到。中位数使\(\sum_i|x_i-c|\)最小：把最小、最大项依次配对，每对的绝对偏差和在\(c\)位于两项之间时最小。偶数时，取中间两项的平均；奇数时，取未配对的中间一项，使其偏差为零。两种选择均能同时使各部分最小。

\subsection{离散程度}

定义描述这组数据的{\heiti 样本方差}与标准差为
\begin{equation}\label{eq:prob-descriptive-variance}
 s_n^2=\frac1n\sum_i(x_i-\bar x)^2
 =\frac1n\sum_ix_i^2-\bar x^2,\qquad s_n=\sqrt{s_n^2}.
\end{equation}
这里分母为\(n\)，是在描述已有数据的平方偏差平均。后面估计未知总体方差时，还会出现分母为\(n-1\)的另一种统计量。

本节数据有\(\bar x=10.5\)、\(s_n^2=4.25\)。若每个值都变为\(y_i=ax_i+b\)，则
\[
 \bar y=a\bar x+b,\qquad s_{y,n}^2=a^2s_{x,n}^2.
\]
还可以用{\heiti 极差}\(x_{(n)}-x_{(1)}\)和{\heiti 四分位距}\(Q_3-Q_1\)描述分散程度。本节二者分别为7和2.5；前者使用两端数据，后者集中于中间一半的位置范围。

\subsection{分组汇总}

设数据分成\(m\)组，第\(h\)组有\(n_h\)个数据，均值为\(\bar x_h\)，描述方差为\(s_h^2\)，总数为\(n=\sum_hn_h\)。直接按组求和，得到
\[
 \bar x=\sum_h\frac{n_h}n\bar x_h.
\]
再写\(x_{hj}-\bar x=(x_{hj}-\bar x_h)+(\bar x_h-\bar x)\)，平方后在各组求和。因为\(\sum_j(x_{hj}-\bar x_h)=0\)，交叉项消失，故
\begin{equation}\label{eq:prob-group-variance}
 s_n^2=\sum_h\frac{n_h}n s_h^2
 +\sum_h\frac{n_h}n(\bar x_h-\bar x)^2.
\end{equation}
合并后的离散程度包含组内波动与组间均值差异。只把各组方差平均，会遗漏第二部分。这与前面条件平均的总方差分解具有相同的有限求和依据。

\begin{example}
甲组20个数据，均值为60、方差为4；乙组30个数据，均值为70、方差为9。求合并后的均值与描述方差，再讨论两组均值相同时公式怎样简化。
\end{example}

\subsection{统计摘要的选择}

仍看式\eqref{eq:prob-description-data}。把最大值15改为30，均值明显上升，而中位数保持不变；极差与平方偏差也对这一变化敏感。若研究总体总量或平方误差，均值和方差很有用；若关注典型位置且存在少量很大的数值，中位数和分位数可以提供另一类信息。

若数据来自不同组别，还需要观察组内分布及其混合比例。同一均值可能来自所有数据集中在中心附近，也可能来自相距很远的两组数据。图形、均值、分位数与离散程度分别保留不同信息，选择摘要时应保留样本量、变量的含义以及数据来源。

\begin{example}[据2025年上海卷改编]\label{ex:prob-swimming-data}
历届奥运会男子\(4\times100\)米混合泳接力冠军的十个成绩数据如下，单位为秒，已按升序排列：
\[
\begin{array}{rrrrr}
206.78&207.46&207.95&209.34&209.35\\
210.68&213.73&214.84&216.93&216.93
\end{array}
\]
求极差、中位数与样本均值；从这十条记录中随机选三条，求恰有两条成绩大于211秒的概率。
\end{example}

\begin{exercise}
\mbox{}\par\nobreak

（1）对数据\(2,3,3,4,5,7,7,9\)，求均值、中位数、众数、四分位距与描述方差，并画箱线图。

（2）（据2020年天津卷第4题改编）由图\ref{fig:prob-histograms}右图，求80个零件中直径在\([5.43,5.47)\)毫米内的个数，并估计该批产品中此类零件的比例。

（3）两组人数相等，均值为\(\mu\pm d\)，方差均为\(v\)。求合并后的均值与方差。

（4）分别构造均值相同而方差不同，以及均值方差相同而形态不同的两组数据。
\end{exercise}

\FloatBarrier
\newpage
\section{抽样分布}
\label{sec:prob-sampling-distributions}

一组数据产生一个均值，另一次抽样可能产生另一个均值。把抽样过程也纳入模型，才能研究这种变化有多大，以及样本量怎样影响估计。这里先给定总体或分布，再研究由它产生的样本。

\subsection{统计量}

\begin{definition}
由样本\((X_1,\ldots,X_n)\)确定、表达式中不含未知参数的函数，称为{\heiti 统计量}。它在抽样模型下的分布称为{\heiti 抽样分布}。
\end{definition}

例如样本均值\(\bar X=n^{-1}\sum_iX_i\)是统计量，抽样前是随机变量；记录到\(x_1,\ldots,x_n\)后，得到观测值\(\bar x\)。样本比例、最大值和样本方差同样是统计量。

从\(\{0,2\}\)中有放回独立抽两次，各值等可能。四条有序样本\((0,0),(0,2),(2,0),(2,2)\)各以\(1/4\)发生，样本均值的分布为
\[
\begin{array}{c|ccc}
\bar X&0&1&2\\\hline P&1/4&1/2&1/4
\end{array}
\]
这是均值本身的分布，而非单次观测的总体分布。

抽样分布也可以借助事件关系求出。若\(X_1,\ldots,X_n\)相互独立，具有同一分布函数\(F\)，样本最大值\(M_n=\max\{X_1,\ldots,X_n\}\)满足
\[
 P(M_n\le t)=P(X_1\le t,\ldots,X_n\le t)=[F(t)]^n.
\]
最大值不超过\(t\)，恰好要求每次观测都不超过\(t\)。这使一个统计量的分布转化为单次观测的概率，无须逐一枚举所有样本。

样本较大时，可以固定一个模型，反复独立生成整组样本，每组计算一次统计量，再整理这些值的频数与图形。这样的模拟近似的是给定模型下的抽样分布；每一组内部包含\(n\)次观测，重复整组的次数则决定模拟记录的数量。

\subsection{样本均值的波动}

下面称具有同一分布且相互独立的\(X_1,\ldots,X_n\)为{\heiti 独立同分布样本}。设总体均值为\(\mu\)、方差为有限的\(\sigma^2\)，由期望线性性与独立时方差相加，
\begin{equation}\label{eq:prob-sample-mean-characteristics}
 E\bar X=\mu,\qquad
 \operatorname{Var}\bar X=\frac1{n^2}\sum_i\operatorname{Var}X_i
 =\frac{\sigma^2}{n}.
\end{equation}
单次观测的标准差为\(\sigma\)，均值的标准差则为\(\sigma/\sqrt n\)，也称样本均值的{\heiti 标准误}。样本量增为四倍，标准误变为原来的一半。

独立二分类试验中，令成功比例为\(\widehat p=n^{-1}\sum_iI_i\)，则
\[
 E\widehat p=p,\qquad \operatorname{Var}\widehat p=\frac{p(1-p)}n.
\]

\begin{figure}[!ht]
\centering
\includegraphics[width=.9\linewidth]{probability-sample-means}
\caption{成功概率为\(0.4\)时，两个样本量下样本比例的精确分布。每根竖线高度为对应取值的概率。}
\end{figure}

有限总体\(z_1,\ldots,z_N\)的均值为\(\mu\)，描述方差为\(\sigma_N^2=N^{-1}\sum_j(z_j-\mu)^2\)。若简单随机不放回抽取\(n\)个，则不同抽取位置的协方差为\(-\sigma_N^2/(N-1)\)。这是因为所有不同下标的中心化乘积之和为\(-\sum_j(z_j-\mu)^2\)。因此
\begin{equation}\label{eq:prob-finite-mean-variance}
 E\bar X=\mu,\qquad
 \operatorname{Var}\bar X=\frac{\sigma_N^2}{n}\frac{N-n}{N-1}\quad(N>1).
\end{equation}
有限总体修正与超几何模型一致；使用时需区分独立重复取样与有限总体不放回取样。

\begin{example}
独立测量的均值为未知常数\(\mu\)，每次测量的标准差为2。用\(n\)次测量的均值估计\(\mu\)。为了使偏离至少1的概率不超过\(0.04\)，根据Chebyshev界应取多大样本量？若精度改为\(0.5\)，样本量怎样变化？
\end{example}

测量误差还可能含有共同的来源。设\(X_i=\mu+Z+\varepsilon_i\)，其中\(Z\)是同一轮测量共有的校准误差，\(\varepsilon_i\)是各次独立误差；它们相互独立、均值为0，方差分别为\(\tau^2\)、\(\sigma^2\)。此时
\[
 \bar X=\mu+Z+\frac1n\sum_i\varepsilon_i,
 \qquad \operatorname{Var}\bar X=\tau^2+\frac{\sigma^2}n.
\]
平均能削弱各次独立误差，却保留共同误差。增加测量次数与改善校准，解决的是不同来源的波动。

\subsection{大数定律}

\begin{theorem}[弱大数定律]\label{thm:prob-wlln}
设\(X_1,X_2,\ldots\)独立同分布，期望为\(\mu\)，方差为有限的\(\sigma^2\)。令\(\bar X_n=n^{-1}\sum_{i=1}^nX_i\)，则对任意\(\varepsilon>0\)，
\begin{equation}\label{eq:prob-wlln}
 P(|\bar X_n-\mu|\ge\varepsilon)\longrightarrow0\qquad(n\to\infty).
\end{equation}
\end{theorem}

\begin{proof}
由式\eqref{eq:prob-sample-mean-characteristics}及Chebyshev不等式，
\[
 0\le P(|\bar X_n-\mu|\ge\varepsilon)
 \le\frac{\sigma^2}{n\varepsilon^2}\longrightarrow0.\qedhere
\]
\end{proof}

这里趋于零的是固定精度以外的偏离概率，称\(\bar X_n\){\heiti 依概率收敛}于\(\mu\)。每次增加样本量，具体观测的误差仍可能变大或变小；定理控制整个抽样模型中偏离发生的可能性。

令\(X_i\)为事件发生的指示变量，就得到频率的稳定依据：
\[
 P(|\widehat p-p|\ge\varepsilon)
 \le\frac{p(1-p)}{n\varepsilon^2}
 \le\frac1{4n\varepsilon^2}.
\]
例如要求偏离不超过\(0.05\)的概率至少为\(0.95\)，采用上述统一界，取\(n\ge2000\)即可。这是由有限样本不等式得到的充分样本量，已知更多分布信息时可以得到更精细的结果。

\begin{exercise}
\mbox{}\par\nobreak

（1）从\(\{0,1,2\}\)等可能有放回抽两次，求样本均值的分布、期望与方差。

（2）总体标准差为6，独立样本量为36及144时，样本均值的标准误分别是多少？用Chebyshev界估计偏离总体均值至少2的概率。

（3）某有限总体有100个数，均值为10、描述方差为9。不放回随机抽25个，求样本均值的期望与方差，并与有放回抽样比较。

（4）用频率估计未知概率，要求偏离小于\(0.02\)的概率至少为\(0.9\)。由统一Chebyshev界给出足够的样本量。

（5）独立掷一枚均匀骰子\(n\)次，记最大点数为\(M_n\)。求\(M_n\)的分布列。若要求至少出现一次6点的概率不小于\(0.99\)，最少应掷多少次？

（6）独立重复掷均匀骰子，点数为\(X_1,\ldots,X_n\)。用大数定律说明，偶数点的频率、\(n^{-1}\sum_iX_i\)及\(n^{-1}\sum_iX_i^2\)分别趋近什么数；分别指出所用的随机变量。

（7）在上述共同误差模型中，设\(\tau^2=1/4\)、\(\sigma^2=4\)。使样本均值的方差不超过\(0.3\)，至少需要多少次测量？若要求标准误不超过\(0.4\)，仅增加测量次数能否做到？
\end{exercise}

\FloatBarrier
\newpage
\subsection{样本方差的修正*}

现在用独立同分布样本估计总体方差。令
\[
 S_n^2=\frac1n\sum_i(X_i-\bar X)^2.
\]
在总体均值\(\mu\)处作平方分解，得到
\[
 \sum_i(X_i-\bar X)^2=\sum_i(X_i-\mu)^2-n(\bar X-\mu)^2.
\]
取期望，由\(E(\bar X-\mu)^2=\sigma^2/n\)，
\[
 E\sum_i(X_i-\bar X)^2=n\sigma^2-\sigma^2=(n-1)\sigma^2.
\]
因此\(ES_n^2=(n-1)\sigma^2/n\)。当\(n\ge2\)时，改用
\begin{equation}\label{eq:prob-unbiased-variance}
 S^2=\frac1{n-1}\sum_i(X_i-\bar X)^2
\end{equation}
便有\(ES^2=\sigma^2\)，称为总体方差的{\heiti 无偏估计量}。

两种分母各有任务：\(n\)描述当前数据的平均平方偏差，\(n-1\)使这一抽样模型下的估计在期望上恢复总体方差。残差满足\(\sum_i(X_i-\bar X)=0\)，落在\(\mathbb R^n\)的一个\(n-1\)维子空间中；前面的期望计算给出了修正倍数的直接依据。

\subsection{正态近似*}

这里先引用正态分布对独立相加的封闭性：若\(X_i\)独立且均为\(N(\mu,\sigma^2)\)，则
\[
 \bar X\sim N\left(\mu,\frac{\sigma^2}n\right),\qquad
 \frac{\sqrt n(\bar X-\mu)}\sigma\sim N(0,1).
\]
这是每个样本量下的精确分布。一般总体并不要求正态，但大样本的标准化均值仍有共同的极限分布。

\begin{theorem}[中心极限定理]
若\(X_i\)独立同分布，期望为\(\mu\)，方差满足\(0<\sigma^2<\infty\)，则对每个实数\(z\)，
\begin{equation}\label{eq:prob-clt}
 P\left(\frac{\sqrt n(\bar X_n-\mu)}\sigma\le z\right)
 \longrightarrow\Phi(z).
\end{equation}
\end{theorem}

大数定律说明平均集中到哪里，中心极限定理则描述按\(\sqrt n\)缩放后的波动。原来的总体可以离散，也可以连续；标准化后的分布函数趋于同一个\(\Phi\)。

下面用逐项替换的方法证明。把和式中的变量依次换成同均值、同方差的正态变量：Taylor展开的一次项与二次项在取期望后相同，只需控制余项。最后还要说明，光滑函数上的比较怎样给出区间概率的比较。

\begin{proof}
令\(Y_i=(X_i-\mu)/\sigma\)，则\(EY_i=0\)、\(EY_i^2=1\)，并记
\[
 T_n=\frac{Y_1+\cdots+Y_n}{\sqrt n}.
\]
另取相互独立的标准正态变量\(Z_1,\ldots,Z_n\)，且使它们与所有\(Y_i\)独立。由前面引用的正态相加性质，
\[
 G_n=\frac{Z_1+\cdots+Z_n}{\sqrt n}\sim N(0,1).
\]
因此只需比较\(T_n\)与\(G_n\)的分布。

\smallskip
{\heiti 1. Taylor余项的统一控制。}
取二阶连续可微的有界函数\(f\)，要求\(f'\)、\(f''\)有界，且\(f''\)在实数轴上一致连续。由积分余项公式，
\[
 f(w+h)=f(w)+f'(w)h+\frac12f''(w)h^2+R(w,h),
\]
其中
\[
 R(w,h)=h^2\int_0^1(1-t)\bigl[f''(w+th)-f''(w)\bigr]\,dt.
\]
记\(M=\sup_x|f''(x)|\)。给定\(\varepsilon>0\)，由一致连续性，可取\(\delta>0\)，使\(|u-v|\le\delta\)时有\(|f''(u)-f''(v)|\le\varepsilon\)。当\(|h|\le\delta\)时，余项不超过\(\varepsilon h^2/2\)；对其他增量，仍有\(|R(w,h)|\le Mh^2\)。合起来便是
\begin{equation}\label{eq:prob-clt-remainder}
 |R(w,h)|\le\frac{\varepsilon}{2}h^2
 +Mh^2\mathbf1_{\{|h|>\delta\}}.
\end{equation}
这个界不依赖展开位置\(w\)，因而可以用于每一次替换。

\smallskip
{\heiti 2. 逐项替换。}
记
\[
 W_{n,k}=\frac{Z_1+\cdots+Z_{k-1}+Y_{k+1}+\cdots+Y_n}{\sqrt n}.
\]
第\(k\)次替换比较的是\(f(W_{n,k}+Y_k/\sqrt n)\)与\(f(W_{n,k}+Z_k/\sqrt n)\)。由于\(W_{n,k}\)分别与\(Y_k,Z_k\)独立，且两者均值都为0、二阶矩都为1，Taylor展开取期望后，常数项、一次项与二次项分别相消。例如
\[
 E\left[f'(W_{n,k})\frac{Y_k}{\sqrt n}\right]=0,\qquad
 E\left[f''(W_{n,k})\frac{Y_k^2}{2n}\right]
 =\frac{E f''(W_{n,k})}{2n},
\]
换成\(Z_k\)所得结果相同。

由式\eqref{eq:prob-clt-remainder}控制两份余项，再将\(n\)次替换的差相加，得到
\begin{align*}
 |Ef(T_n)-Ef(G_n)|
 &\le\varepsilon
 +M E\bigl(Y_1^2\mathbf1_{\{|Y_1|>\delta\sqrt n\}}\bigr)\\
 &\quad+M E\bigl(Z_1^2\mathbf1_{\{|Z_1|>\delta\sqrt n\}}\bigr).
\end{align*}
若\(EV^2<\infty\)，则\(E(V^2\mathbf1_{\{|V|>r\}})\to0\)：它就是二阶矩对应的非负求和或积分在\([-r,r]\)之外的尾部，随区间扩大而趋于零。因此固定上述\(\delta\)，令\(n\to\infty\)，右端两项尾部都趋于零。再由\(\varepsilon\)任意，得到
\begin{equation}\label{eq:prob-clt-smooth}
 Ef(T_n)\longrightarrow Ef(Z),\qquad Z\sim N(0,1).
\end{equation}
这里\(G_n\)与\(Z\)具有相同分布，所以\(Ef(G_n)=Ef(Z)\)。

\smallskip
{\heiti 3. 从光滑函数回到分布函数。}
区间\(({-\infty},z]\)的指示函数在\(z\)处跳跃，不能直接代入上面的Taylor展开。我们用两个光滑函数从上下夹住它。令
\[
 H(u)=\begin{cases}
 1,&u\le0,\\
 1-10u^3+15u^4-6u^5,&0<u<1,\\
 0,&u\ge1.
 \end{cases}
\]
在\([0,1]\)内，\(H'(u)=-30u^2(1-u)^2\le0\)，且一、二阶导数在两端均为0。因此\(H\)二阶连续可微，满足前面对\(f\)的全部要求。

给定\(a>0\)，令
\[
 f_-(t)=H\left(\frac{t-z+a}{a}\right),\qquad
 f_+(t)=H\left(\frac{t-z}{a}\right).
\]
前者在\([z-a,z]\)内从1降到0，后者在\([z,z+a]\)内从1降到0，故
\[
 f_-(t)\le\mathbf1_{\{t\le z\}}\le f_+(t).
\]
对标准正态变量\(Z\)，还有\(\Phi(z-a)\le Ef_-(Z)\)、\(Ef_+(Z)\le\Phi(z+a)\)。对\(f_-\)、\(f_+\)分别使用式\eqref{eq:prob-clt-smooth}，可知任给\(\eta>0\)，当\(n\)充分大时，
\[
 \Phi(z-a)-\eta\le P(T_n\le z)\le\Phi(z+a)+\eta.
\]
由\(\Phi\)的连续性，先取充分小的\(a\)，再取充分小的\(\eta\)，便得\(P(T_n\le z)\to\Phi(z)\)。
\end{proof}

这也解释了正态近似的一个共同来源：大量独立小项相加时，均值与方差决定Taylor展开的前两项，其余差异在求和与标准化后逐渐消失。

若\(K\sim B(n,p)\)，它是Bernoulli变量之和，因此可用均值\(np\)、方差\(np(1-p)\)的正态模型近似。整数事件\(a\le K\le b\)通常用连续区间\([a-1/2,b+1/2]\)近似，这称为{\heiti 连续性修正}。例如\(K\sim B(100,0.4)\)时，
\[
 P(35\le K\le45)\approx
 \Phi\left(\frac{45.5-40}{\sqrt{24}}\right)
 -\Phi\left(\frac{34.5-40}{\sqrt{24}}\right).
\]
近似的质量还受\(p\)及尾部位置影响。对具体问题，可以把二项精确概率与近似值比较；样本量本身不是精度的完整说明。

\begin{exercise}
\mbox{}\par\nobreak

（1）独立样本量为5，观察到\(1,2,2,3,7\)。分别计算描述方差与无偏样本方差，并说明各自的用途。

（2）从\(N(10,4)\)中独立抽取16个观测，求样本均值在9至11之间的概率。这里是否需要使用中心极限定理？

（3）\(K\sim B(200,0.5)\)，写出\(90\le K\le110\)的精确概率与连续性修正后的正态近似表达，并用计算工具比较数值。

（4）独立取\(X_1,\ldots,X_{100}\sim U(0,1)\)。用中心极限定理近似计算\(P(0.45\le\bar X\le0.55)\)，再估计总和超过55的概率。

（5）设\(X_i\)独立同分布，均值为\(\mu\)，方差为\(\sigma^2>0\)。分别考察固定\(c>0\)时的\(P(|\bar X_n-\mu|\le c)\)与\(P(|\bar X_n-\mu|\le c\sigma/\sqrt n)\)：它们的极限分别是什么？说明大数定律与中心极限定理在这里各自描述了什么。

（6）沿用证明中的\(Y_i,Z_i,T_n,G_n\)。进一步设\(E|Y_1|^3<\infty\)，且有界函数\(f\)三阶连续可微，各阶导数有界，\(|f'''|\le L\)。直接使用三阶余项，证明
\[
 |Ef(T_n)-Ef(G_n)|\le
 \frac{L}{6\sqrt n}\bigl(E|Y_1|^3+E|Z_1|^3\bigr).
\]
比较这个估计与正文的余项控制，说明正文为何只需有限二阶矩。
\end{exercise}

\FloatBarrier
\newpage
\section{参数估计}
\label{sec:prob-estimation}

前一节给定总体，研究样本会怎样变化。现在总体的某些参数未知，已经取得样本，希望由它估计参数。一个估计方法需要同时说明计算规则，以及规则在抽样模型中的表现。

\subsection{点估计}

\begin{definition}
用统计量\(T=T(X_1,\ldots,X_n)\)估计参数\(\theta\)，称\(T\)为\(\theta\)的{\heiti 估计量}。代入实际观测所得的数称为{\heiti 估计值}。
\end{definition}

独立同分布样本中，\(\bar X\)用于估计总体均值\(\mu\)，成功比例\(\widehat p\)用于估计二分类概率\(p\)。大数定律说明，在已建立的条件下，样本量增大时它们偏离相应参数固定幅度的概率趋零。这为使用平均与频率提供了依据。

一种自然出发点是用样本平均匹配理论平均。例如\(X\sim U(0,\theta)\)时，理论均值为\(\theta/2\)，把它与\(\bar X\)对应，得到\(2\bar X\)。这种由理论矩与样本矩建立参数方程的想法称为{\heiti 矩估计}。所得规则的质量仍需要在具体模型下检验。

正态模型中的均值、中位数和对称中心相同，可用样本均值或样本中位数估计其位置。同一参数因而可以有不同估计规则，评价时需结合模型、误差目标和取样方式。

\begin{example}
从\(U(0,\theta)\)独立抽取\(n\)个观测，用\(2\bar X\)估计\(\theta\)。求这个估计量的期望与方差，并由Chebyshev界说明它随样本量增大趋近\(\theta\)的依据。
\end{example}

这个模型还给出另一条估计思路：总体端点限制了所有观测，样本最大值\(M_n=\max_iX_i\)可以估计\(\theta\)。由上一节最大值的分布公式，对\(0<\varepsilon<\theta\)，
\[
 P(|M_n-\theta|\ge\varepsilon)
 =P(M_n\le\theta-\varepsilon)
 =\left(1-\frac\varepsilon\theta\right)^n\longrightarrow0.
\]
最大值利用了已知的端点信息，虽不超过\(\theta\)，仍能随样本量增加而逼近\(\theta\)。

\subsection{估计误差}

\begin{definition}
估计量\(T\)的{\heiti 偏差}为\(E_\theta T-\theta\)，{\heiti 均方误差}为
\[
 \operatorname{MSE}_\theta(T)=E_\theta(T-\theta)^2.
\]
若对所考察模型中的参数值均有\(E_\theta T=\theta\)，称\(T\)为{\heiti 无偏估计量}。下标表示期望按参数为\(\theta\)的模型计算。
\end{definition}

由平方误差分解，
\begin{equation}\label{eq:prob-bias-variance}
 \operatorname{MSE}_\theta(T)=\operatorname{Var}_\theta T
 +(E_\theta T-\theta)^2.
\end{equation}
无偏性约束平均位置，均方误差还计入抽样波动。例如\(\bar X\)无偏估计\(\mu\)，其均方误差为\(\sigma^2/n\)；一个有偏规则也可能在某些参数范围内具有更小的均方误差。

\begin{example}
设\(K\sim B(n,p)\)，比较估计量\(T_1=K/n\)与\(T_2=(K+1)/(n+2)\)。求各自的偏差和均方误差。取\(n=10\)，分别在\(p=0.5\)与\(p=0.1\)时比较，判断是否能仅凭无偏性决定优劣。
\end{example}

后一估计可写成\(T_2=\dfrac n{n+2}\widehat p+\dfrac2{n+2}\cdot\dfrac12\)，即把样本比例向\(1/2\)略微移动。一般地，给定参考值\(m\)，取\(0\le c\le1\)，令\(T_c=m+c(\bar X-m)\)，则
\[
 \operatorname{MSE}_\mu(T_c)
 =c^2\frac{\sigma^2}n+(1-c)^2(\mu-m)^2.
\]
减小\(c\)会减小方差，同时增加偏离参考值时的平方偏差。哪个影响更大，取决于总体均值的位置以及样本量。

这些误差是在给定抽样机制下计算的。若取样来源与目标总体不同，或测量含系统偏差，应先修正模型和数据获取；单纯增加样本量只能按原模型减小相应的随机波动。

\begin{exercise}
\mbox{}\par\nobreak

（1）估计量\(T\)的偏差为\(0.2\)，方差为\(0.04\)，求均方误差。另一无偏估计量的方差为\(0.1\)，按均方误差比较二者。

（2）从均值\(\mu\)、方差\(\sigma^2\)的模型中取独立样本。比较只用\(X_1\)与使用\(\bar X\)估计\(\mu\)时的偏差与均方误差。

（3）独立抽检100件产品，发现6件不合格，以样本比例估计不合格率。方案甲逐件检验，检验费2元，不合格品另付更换费10元；方案乙免检出售，每件不合格品赔付50元。比较估计的逐件期望费用，并求两方案期望费用相同的临界不合格率。

（4）从\(U(0,\theta)\)独立抽样，使样本最大值距\(\theta\)不足\(0.1\theta\)的概率至少为\(0.95\)，最少需多少个观测？若只用\(X_1\)估计端点，增加样本总量能否减小误差？

（5）设\(E\bar X=\mu\)、\(\operatorname{Var}\bar X=1\)。用\(T=\bar X/2\)估计\(\mu\)，求使它的均方误差小于\(\bar X\)的全部\(\mu\)值。说明这种比较为什么依赖参数所在的范围。

（6）独立同分布样本的均值为\(\mu\)、方差为\(\sigma^2\)。令\(T_n=\bar X+1/\sqrt n\)，求其偏差、均方误差，证明它仍依概率收敛于\(\mu\)。
\end{exercise}

\FloatBarrier
\newpage
\subsection{似然*}

设独立二分类试验进行\(n\)次，成功概率\(p\)未知。观察到一条确定的0、1序列\(x_1,\ldots,x_n\)，其中成功\(k\)次。不同的\(p\)给这条记录分配的概率为
\begin{equation}\label{eq:prob-likelihood}
 L(p)=p^k(1-p)^{n-k},\qquad 0\le p\le1.
\end{equation}
固定数据，把它的模型概率看作参数的函数，称为{\heiti 似然函数}；使其最大的参数称为{\heiti 极大似然估计}。似然比较不同模型对当前记录的解释，不给参数值本身分配概率。

当\(0<k<n\)时，取对数不改变最大值位置，并有
\[
 \ell(p)=k\ln p+(n-k)\ln(1-p),\qquad
 \ell'(p)=\frac{k-np}{p(1-p)}.
\]
导数在\(p=k/n\)前为正、后为负，故最大值在\(\widehat p=k/n\)。全失败或全成功时分别取\(\widehat p=0,1\)。若只记录成功次数，似然多出与\(p\)无关的因子\(\binom nk\)，最大值位置仍然相同。

样本比例由匹配平均得到，也由当前模型的似然最大化得到。两个出发点在这个模型中给出同一个规则，而它的偏差与波动仍由抽样分布研究。

\begin{exercise}
独立试验20次，观察到7次成功。写出似然函数，求极大似然估计，并比较\(p=0.2\)、\(p=0.35\)、\(p=0.5\)时的似然。说明这里哪一部分是已固定的数据。
\end{exercise}

\subsection{置信区间*}

一个估计值本身没有说明抽样误差。可以让样本产生一个区间，并研究这一构造在反复抽样时覆盖真参数的概率。

\begin{definition}
给定\(0<\alpha<1\)。若由样本确定的区间\(I(X_1,\ldots,X_n)\)，对模型中的参数满足
\[
 P_\theta\{\theta\in I(X_1,\ldots,X_n)\}\ge1-\alpha,
\]
则称其为覆盖率至少\(1-\alpha\)的{\heiti 置信区间}，\(1-\alpha\)称为置信水平。这里参数固定，区间随样本变化。
\end{definition}

先设独立同分布样本的总体方差不超过已知数\(v\)。Chebyshev界给出
\[
 P\left(|\bar X-\mu|\ge\sqrt{\frac{v}{n\alpha}}\right)\le\alpha,
\]
可取区间
\begin{equation}\label{eq:prob-chebyshev-interval}
 \left[\bar X-\sqrt{\frac{v}{n\alpha}},\,
       \bar X+\sqrt{\frac{v}{n\alpha}}\right]
\end{equation}
其覆盖率至少为\(1-\alpha\)。它只利用方差上界，因而通常较保守。

若进一步知道总体为\(N(\mu,\sigma^2)\)，且\(\sigma\)已知，则标准化样本均值精确服从\(N(0,1)\)。令\(q_\alpha>0\)满足\(\Phi(q_\alpha)=1-\alpha/2\)，有
\[
 P\left(-q_\alpha\le\frac{\sqrt n(\bar X-\mu)}\sigma\le q_\alpha\right)=1-\alpha.
\]
等价改写事件，得到覆盖率为\(1-\alpha\)的区间
\begin{equation}\label{eq:prob-normal-interval}
 \left[\bar X-q_\alpha\frac\sigma{\sqrt n},\,
       \bar X+q_\alpha\frac\sigma{\sqrt n}\right].
\end{equation}
95\%置信水平下\(q_\alpha\approx1.96\)。模型信息增加，使构造可以利用更精确的分布；两种区间的依据分别是概率界与精确抽样分布。

\begin{figure}[!ht]
\centering
\includegraphics[width=.55\linewidth]{probability-confidence}
\caption{从\(N(0,1)\)反复独立抽取25个观测所得的12个模拟区间。竖线是真均值，区间随样本移动。}
\end{figure}

观测以后，区间已经确定，覆盖或不覆盖也是确定事实；置信水平描述的是整个抽样程序。增大\(n\)可缩短区间，要求更高覆盖率则通常扩大区间。

\begin{exercise}
\mbox{}\par\nobreak

（1）独立正态样本量为100，总体标准差已知为2，观测均值为12.3。分别写出95\%正态置信区间及由方差界得到的95\%保守区间。

（2）在已知标准差的正态模型下，要将区间长度减半，样本量应怎样改变？保持样本量不变，把置信水平从95\%提高到99\%，区间怎样变化？

（3）解释“这套区间构造的覆盖率为95\%”与“这一次取得的区间包含未知参数”分别是哪一种陈述。

（4）设\(K\sim B(n,p)\)、\(\widehat p=K/n\)。由\(\operatorname{Var}\widehat p\le1/(4n)\)构造覆盖率至少为\(1-\alpha\)的\(p\)区间。将该区间与\([0,1]\)取交，是否会降低覆盖率？
\end{exercise}

\FloatBarrier
\newpage
\section{假设检验}
\label{sec:prob-testing}

估计直接寻找参数的可能位置。另一类任务先提出一个模型假设，再比较观测结果与该假设下的随机波动。如果预先选定的偏离在这个模型下很少发生，就据此作出拒绝假设的判断。

\subsection{假设与随机波动}

设独立抛一枚硬币10次，\(K\)为正面次数，未知正面概率为\(p\)。为研究硬币是否偏向正面，提出
\[
 H_0:p=\frac12,\qquad H_1:p>\frac12.
\]
\(H_0\)称为{\heiti 原假设}，\(H_1\)为所考察的{\heiti 备择假设}。在\(H_0\)下，\(K\sim B(10,1/2)\)，因此可以计算各种观测的概率。

偏向正面的证据由较大的\(K\)提供。开始观察前，应确定用哪个统计量、按哪个方向偏离以及怎样作出判断。观察后再根据结果选择方向，会改变原来概率计算所对应的程序。

假设检验始终需要保留随机波动的可能性。即使\(H_0\)成立，也可能出现很大的\(K\)；下面通过拒绝规则控制这种错误发生的概率。

\subsection{拒绝域}

\begin{definition}
若统计量落入预先指定的集合\(R\)就拒绝\(H_0\)，称\(R\)为{\heiti 拒绝域}。若\(H_0\)成立时拒绝的概率不超过\(\alpha\)，称此规则的显著性水平不超过\(\alpha\)。
\end{definition}

在前面的硬币模型中，选择\(R=\{K\ge9\}\)，则
\[
 P_{1/2}(K\ge9)=\frac{\binom{10}9+\binom{10}{10}}{2^{10}}
 =\frac{11}{1024}\approx0.0107.
\]
若扩大为\(K\ge8\)，概率变为\(56/1024\approx0.0547\)。因此要求水平不超过\(0.05\)时，这类整数阈值规则可选择\(K\ge9\)，实际错误概率小于给定上限。

观测到\(K=9\)，便按既定规则拒绝公平假设，认为数据提供了偏向正面的证据。观测到\(K=7\)时不拒绝，表示当前数据在这一规则下尚不足以支持拒绝；这不等于证明\(p=1/2\)。

若改问硬币是否公平，备择为\(H_1:p\ne1/2\)，正面次数过大或过小都构成偏离。取双侧拒绝域\(K\le1\)或\(K\ge9\)，原假设下两端概率相加为\(22/1024\approx0.0215\)。

\begin{example}
某工艺原不合格率为\(0.1\)。为判断不合格率是否升高，独立抽检20件，预先规定发现至少5件不合格就拒绝原假设。写出原假设、备择、统计量和拒绝域，计算原假设成立时错误拒绝的概率。
\end{example}

\subsection{列联表}

两个分类变量的观测频数可以整理为\(2\times2\)列联表：
\[
\begin{array}{c|cc|c}
 & B&B^c&\text{合计}\\\hline
 A&a&b&a+b\\
 A^c&c&d&c+d\\\hline
 \text{合计}&a+c&b+d&N
\end{array}
\]
现在考察总体中两分类变量是否独立。原假设为\(P(A\cap B)=P(A)P(B)\)。用样本边际比例估计两类概率，独立假设下第一格的期望频数为\((a+b)(a+c)/N\)，其余格同理由行合计乘列合计再除以\(N\)得到。

记观测频数为\(O_{ij}\)，拟合的期望频数为\(E_{ij}\)。{\heiti Pearson统计量}定义为
\begin{equation}\label{eq:prob-pearson}
 \chi^2=\sum_{i,j}\frac{(O_{ij}-E_{ij})^2}{E_{ij}}
 =\frac{N(ad-bc)^2}{(a+b)(c+d)(a+c)(b+d)}.
\end{equation}
用样本边际拟合后，观测与期望的各行、各列之差都合计为0。因此只要知道第一格的差\(\delta\)，其余三格也随之确定：
\[
 \delta=a-\frac{(a+b)(a+c)}N=\frac{ad-bc}N,
 \qquad (O_{ij}-E_{ij})=
 \begin{pmatrix}\delta&-\delta\\-\delta&\delta\end{pmatrix}.
\]
把四个平方差代回统计量，得到\(\chi^2=\delta^2\sum_{i,j}1/E_{ij}\)，通分便得式\eqref{eq:prob-pearson}。\(ad=bc\)时，四格都恰好等于拟合期望；偏离越大，统计量越大。

标准正态变量\(Z\)的平方\(Z^2\)所服从的分布，称为自由度为1的\(\chi^2\)分布。在独立观测、各格期望频数足够大的抽样模型中，原假设下统计量\(\chi^2\)近似服从此分布。本章引用这一近似结论及其临界概率，使用时要求各格拟合期望频数至少为5。常用表为
\[
\begin{array}{c|ccc}
 \alpha&0.05&0.01&0.001\\\hline
 \chi^2_\alpha&3.841&6.635&10.828
\end{array}
\]
其含义是原假设下\(P(\chi^2\ge\chi^2_\alpha)\approx\alpha\)。据此在\(\chi^2\ge\chi^2_\alpha\)时拒绝独立假设，结论的水平由这个假设下的尾概率解释。

样本量也影响这一数值。若另一次调查的四格频数恰为原表的\(m\)倍，各组比例保持不变，而式\eqref{eq:prob-pearson}给出新统计量为\(m\chi^2\)。同样的比例差异，在较大的独立样本中可以提供更强的拒绝证据；统计量本身不能直接充当关联强弱的尺度。

\begin{example}[据2025年全国I卷改编]\label{ex:prob-medical-table}
从做过超声波检查的人群中随机调查1000人，得到下表。
\[
\begin{array}{c|rr|r}
 &\text{正常}&\text{不正常}&\text{合计}\\\hline
 \text{患该疾病}&20&180&200\\
 \text{未患该疾病}&780&20&800\\\hline
 \text{合计}&800&200&1000
\end{array}
\]
用频率估计检查结果不正常者患该疾病的概率。再根据\(\alpha=0.001\)的独立性检验，分析检查结果与是否患该疾病的关系，说明判断对应的人群范围。
\end{example}

\begin{exercise}
\mbox{}\par\nobreak

（1）在10次抛币模型中，分别以\(K\ge8\)、\(K\ge9\)为拒绝域，比较原假设下的错误拒绝概率。

（2）两种教学安排各观察50人，甲组达标40人、未达标10人，乙组达标30人、未达标20人。写出列联表及期望频数，在\(\alpha=0.05\)水平作独立性检验。

（3）说明为什么“在1\%水平拒绝独立假设”不能解释为“两个变量独立的概率只有1\%”。

（4）某调查仅记录自愿填写网络问卷的人。即使计算得到很大的Pearson统计量，把结论用于目标总体以前，还需检查什么数据获取条件？

（5）把第（2）题四格频数全部加倍，比较各组达标比例及Pearson统计量。在\(\alpha=0.01\)水平，原表与新表的检验结论是否相同？

（6）由\(Z^2\)的分布及正态对称性，说明自由度为1的\(\chi^2\)分布满足\(P(Z^2\ge c)=2[1-\Phi(\sqrt c)]\)。用正态概率表说明前两个临界值的来源。

（7）两组各\(n\)人，甲、乙达标比例为\(60\%\)、\(40\%\)，\(n\)为5的倍数。求\(\chi^2\)关于\(n\)的表达式，并求在\(0.05\)水平拒绝独立假设的最小\(n\)，核对期望频数条件。
\end{exercise}

\FloatBarrier
\newpage
\subsection{检验误差*}

原假设成立时错误拒绝，称为{\heiti 第一类错误}；某个备择模型成立时未能拒绝原假设，称为{\heiti 第二类错误}。在给定备择参数\(p\)下，拒绝概率称为该处的{\heiti 检验效能}，等于1减去第二类错误概率。

对\(K\ge9\)的硬币规则，若实际\(p=0.8\)，效能为
\[
 P_{0.8}(K\ge9)=\binom{10}9(0.8)^9(0.2)+(0.8)^{10},
\]
第二类错误概率则为其补数。扩大拒绝域会提高发现这一偏离的机会，同时也增加原假设成立时的错误拒绝。样本量和所考察备择的差异，共同影响两种错误的取舍。

在规则\(K\ge c\)下，把拒绝概率写成参数的函数，便得到
\[
 \pi_{n,c}(p)=\sum_{k=c}^n\binom nkp^k(1-p)^{n-k}.
\]
同一个函数在\(p=1/2\)处给出错误拒绝概率，在\(p>1/2\)处给出发现相应偏离的概率。例如比较以下两条事先选定的规则：
\[
 \begin{array}{c|cc}
 \text{拒绝规则}&p=1/2\text{时拒绝的概率}&p=0.8\text{时拒绝的概率}\\\hline
 n=10,\ K\ge9&0.0107&0.3758\\
 n=20,\ K\ge15&0.0207&0.8042
 \end{array}
\]
两者都把原假设下的错误概率控制在\(0.05\)以内，后一规则发现\(p=0.8\)这一偏离的机会更大。小的显著性水平只约束原假设下的表现；能否发现所关心的偏离，还要考察备择下的抽样分布。

对预先确定的单侧大值偏离，若观测到\(K=k\)，定义{\heiti p值}为\(P_{H_0}(K\ge k)\)。它计算原假设下至少同样偏大的结果的概率。在10次抛币模型中，观察到9次正面的p值为\(11/1024\)，观察到8次则为\(56/1024\)。按给定水平比较p值与\(\alpha\)，与对应的尾部拒绝域规则一致。

\subsection{条件检验*}

继续研究列联表。固定两行的对象数，并在独立假设下把第二类变量视为共同成功概率\(p\)的独立0、1记录。给定全部成功数后，任何具有该成功总数的标记方式都有相同概率因子\(p^M(1-p)^{N-M}\)，故条件下这些标记方式等可能。

记总数为\(N\)、成功总数为\(M\)，第一行有\(n\)个对象，其成功数\(X\)的条件分布为
\begin{equation}\label{eq:prob-fisher-conditional}
 \begin{aligned}
 P(X=k\mid\text{边际合计})
 &=\frac{\binom nk\binom{N-n}{M-k}}{\binom NM}\\
 &=\frac{\binom Mk\binom{N-M}{n-k}}{\binom Nn}.
 \end{aligned}
\end{equation}
未知的共同成功概率从条件分布中消去。若事先研究第一行成功比例偏高的方向，就用\(X\ge x_{\rm obs}\)的条件尾概率作判断。

例如两组各6个对象，总成功数6个，第一组成功5个，单侧尾概率为
\[
 \frac{\binom65\binom61+\binom66\binom60}{\binom{12}6}
 =\frac{37}{924}\approx0.0400.
\]
它直接由有限条件模型得到，不需要大频数近似。对例\ref{ex:prob-medical-table}，若按异常检查结果为成功，则固定边际后的\(N=1000,M=200,n=200\)，观测值为180；同一公式给出相应的条件尾概率表达。

边际固定时，列联表只有第一格\(X\)可以变化，其在独立模型下的条件期望为\(nM/N\)。将四格写成\(X,n-X,M-X,N-n-M+X\)，代入式\eqref{eq:prob-pearson}，得到
\[
 \chi^2=\frac{N^3}{n(N-n)M(N-M)}
 \left(X-\frac{nM}N\right)^2.
\]
因此两种方法都研究这一格偏离独立模型中心的程度。条件检验直接使用超几何尾概率，Pearson检验则使用前面引用的大频数近似；二者在同一张表上有共同的研究对象。

\begin{exercise}
\mbox{}\par\nobreak

（1）对硬币模型取备择\(p=0.7\)，计算\(K\ge9\)规则的效能，与\(p=0.8\)比较。

（2）两组各5个对象，总成功数4个，第一组成功4个。写出第一组成功数的条件分布，求单侧尾概率。

（3）在已知方差的正态均值模型中，证明双侧置信区间不包含\(\mu_0\)，等价于标准化样本均值落在相应的双侧尾部。由此构造检验\(H_0:\mu=\mu_0\)的规则。

（4）对\(1\le c\le n\)，通过求导与合并相邻项，证明\(\pi'_{n,c}(p)=n\binom{n-1}{c-1}p^{c-1}(1-p)^{n-c}\)。这说明成功概率升高时，单侧规则的效能怎样变化？

（5）两组各6个对象、总成功数为6时，第一组成功5个与成功1个分别对应哪个方向的偏离？若事先考察两组成功概率是否不同，采用\(X\le1\)或\(X\ge5\)的拒绝域，求其条件概率。
\end{exercise}

\FloatBarrier
\newpage
\section{回归模型}
\label{sec:prob-regression}

当每个对象同时记录两列数值，需要保留它们之间的配对。我们既可以描述共同变化，也可以用一条简单曲线近似两者关系，再根据模型进行预测。

\subsection{成对数据}

设数据为\((x_i,y_i)\)，\(i=1,\ldots,n\)。每对对应同一个观测单位，\(x\)为解释变量，\(y\)为响应变量。散点图将各对数据画成平面上的点，显示趋势、分散程度及异常位置。

以下使用教学数据：五个对象的练习时间\(x\)与测验得分\(y\)为
\begin{equation}\label{eq:prob-regression-data}
\begin{array}{c|rrrrr}
 x\ (\text{小时})&1&2&3&4&5\\\hline
 y\ (\text{分})&20&40&50&40&50
\end{array}
\end{equation}
这组数据用于说明计算与解释，不代表一次实际调查。分别重排两行会改变配对，也就改变所研究的共同变化。

记
\[
 S_{xx}=\sum_i(x_i-\bar x)^2,\quad
 S_{yy}=\sum_i(y_i-\bar y)^2,\quad
 S_{xy}=\sum_i(x_i-\bar x)(y_i-\bar y).
\]
当两列都不恒定时，定义{\heiti 样本相关系数}
\begin{equation}\label{eq:prob-sample-correlation}
 r=\frac{S_{xy}}{\sqrt{S_{xx}S_{yy}}}.
\end{equation}
它与前面按概率加权的相关系数形式一致，这里把已有的\(n\)对记录等权整理。

\subsection{相关的向量解释}

把两列数据中心化，得到\(\mathbb R^n\)中的向量
\[
 \boldsymbol x_c=(x_1-\bar x,\ldots,x_n-\bar x),\qquad
 \boldsymbol y_c=(y_1-\bar y,\ldots,y_n-\bar y).
\]
则\(S_{xy}=\boldsymbol x_c\cdot\boldsymbol y_c\)，\(S_{xx}=\|\boldsymbol x_c\|^2\)，\(S_{yy}=\|\boldsymbol y_c\|^2\)，故\(r\)是两向量夹角的余弦。这直接给出\(|r|\le1\)，并说明正、负号及线性关系的意义。

对教学数据，\(\bar x=3,\bar y=40\)，且
\[
 S_{xx}=10,\qquad S_{yy}=600,\qquad S_{xy}=60,
\]
所以\(r=\sqrt{0.6}\)，描述这五对记录的线性共同变化。

\subsection{最小二乘直线}

用直线\(\widehat y=a+bx\)近似散点，观测与拟合之差\(e_i=y_i-a-bx_i\)称为{\heiti 残差}。选择使残差平方和
\[
 Q(a,b)=\sum_i(y_i-a-bx_i)^2
\]
最小的直线，称为{\heiti 最小二乘回归直线}。

将各列分成均值与中心化部分，交叉项求和为零，得到
\[
 Q(a,b)=n(\bar y-a-b\bar x)^2+S_{yy}-2bS_{xy}+b^2S_{xx}.
\]
当\(S_{xx}>0\)时，再配方，最小值恰在
\begin{equation}\label{eq:prob-least-squares}
 \widehat b=\frac{S_{xy}}{S_{xx}},\qquad
 \widehat a=\bar y-\widehat b\bar x
\end{equation}
处取得。这一推导只需要有限平方展开。

回归直线经过\((\bar x,\bar y)\)，残差满足\(\sum_ie_i=0\)与\(\sum_ix_ie_i=0\)。对教学数据，回归直线为\(\widehat y=22+6x\)，残差依次为\(-8,6,10,-6,-2\)，平方和为240。

\begin{figure}[!ht]
\centering
\includegraphics[width=.88\linewidth]{probability-regression}
\caption{左图为同一组成对数据及最小二乘直线，右图为残差随解释变量的变化。}
\end{figure}

残差图保留直线未解释的部分。若残差有明显弯曲趋势，或分散程度随\(x\)系统改变，使用同一线性模型作预测时应重新检查相应假设。

\subsection{预测的解释}

对给定\(x_0\)，\(\widehat a+\widehat b x_0\)给出模型拟合或预测值。把它解释为同类对象的平均趋势，需要假定新对象与原数据遵循相应关系；单个对象还会有未被直线解释的波动。远离已有数据范围的外推，增加了对关系继续保持的依赖。

可以用一个简单的随机模型明确这种区别。固定\(x_1,\ldots,x_n\)，设\(Y_i=a+bx_i+\varepsilon_i\)，各误差相互独立，均值为0、方差为\(\sigma^2\)。代入最小二乘公式，
\[
 \widehat b=b+\frac{\sum_i(x_i-\bar x)\varepsilon_i}{S_{xx}},
 \qquad E\widehat b=b,\quad
 \operatorname{Var}\widehat b=\frac{\sigma^2}{S_{xx}}.
\]
斜率估计的波动由误差大小和解释变量的分散程度共同决定。预测同类新对象时，拟合值估计平均响应\(a+bx_0\)，而单个新观测还含有自身的随机误差。

相关与回归首先描述观测关系。若练习时间与得分同时受到原有基础、教学安排等因素影响，当前散点图本身不能确定改变练习时间会造成怎样的得分变化。因果问题需回到研究设计。

对正值响应，有时考虑指数趋势，先拟合\(\ln y\)关于\(x\)的直线，再取指数恢复预测。此时最小化的是
\[
 \sum_i(\ln y_i-a-bx_i)^2,
\]
它与最小化原尺度上的\(\sum_i(y_i-e^{a+bx_i})^2\)是不同的误差任务。选择变换时，应同时说明模型关系与所用的误差度量。

\begin{example}[据2025年上海卷改编]
沿用例\ref{ex:prob-swimming-data}的十个游泳成绩，已知成绩\(y\)关于年份\(x\)的最小二乘回归直线为\(y=-0.311x+\widehat b\)，年份的样本均值为2006。求\(\widehat b\)，并预测2028年的成绩，精确到0.01秒。说明此预测与观测数据、趋势延续假设之间的关系。
\end{example}

\begin{exercise}
\mbox{}\par\nobreak

（1）对数据\((1,2),(2,3),(3,5),(4,4)\)，求均值、样本相关系数与最小二乘直线，列出残差并核对残差和。

（2）（据2026年上海卷第17题改编）颗粒物密度关于年份\(x\)的拟合式为\(\widehat y_1=106.544e^{-0.461(x-2014)}\)、\(\widehat y_2=a(x-2014)+83.743\)，其中\(\bar x=2018\)、\(\bar y=301.49/9\)。求\(a\)及两模型对2023年实测值3.88的绝对误差（保留两位小数）。可否据此决定今后的模型？

（3）某正值数据满足\(y_i=C e^{kx_i}\)，且\(x_i\)不全相同。证明对数线性拟合能恢复\(\ln C,k\)。若各\(y_i\)含不同误差，为什么两种平方误差任务通常给出不同结果？

（4）两列数据均不恒定。分别拟合\(y\)关于\(x\)及\(x\)关于\(y\)的最小二乘直线，证明两条直线的斜率乘积为\(r^2\)。当\(0<|r|<1\)时，后一拟合能否由前一条直线直接反解得到？

（5）在固定解释变量的随机模型中取\(\sigma^2=4\)，使用教学数据的五个\(x\)值，求斜率估计的标准差。若在每个\(x\)值处各独立观测四次，新斜率估计的标准差是多少？
\end{exercise}

\FloatBarrier
\newpage
\subsection{回归的投影结构*}

记\(\boldsymbol1=(1,\ldots,1)\)，将两列数据视为向量\(\boldsymbol x,\boldsymbol y\)。全体拟合向量为
\[
 W=\operatorname{span}\{\boldsymbol1,\boldsymbol x\},\qquad
 \widehat{\boldsymbol y}=a\boldsymbol1+b\boldsymbol x.
\]
当\(x\)不恒定时，\(W\)为二维子空间。残差条件\(\sum_ie_i=\sum_ix_ie_i=0\)恰好表示\(\boldsymbol e\perp W\)。因此最小二乘拟合是\(\boldsymbol y\)在\(W\)上的正交投影。对任意\(\boldsymbol v\in W\)，
\[
 \|\boldsymbol y-\boldsymbol v\|^2
 =\|\boldsymbol e\|^2+\|\widehat{\boldsymbol y}-\boldsymbol v\|^2,
\]
从而投影的确使平方误差最小。

又因\(\widehat{\boldsymbol y}-\bar y\boldsymbol1\in W\)，
\[
 \boldsymbol y-\bar y\boldsymbol1
 =(\widehat{\boldsymbol y}-\bar y\boldsymbol1)+\boldsymbol e
\]
是正交分解，给出{\heiti 变差分解}
\begin{equation}\label{eq:prob-regression-decomposition}
 \sum_i(y_i-\bar y)^2=\sum_i(\widehat y_i-\bar y)^2+\sum_ie_i^2.
\end{equation}
当\(S_{yy}>0\)时，定义拟合所解释的平方变差比例
\[
 R^2=1-\frac{\sum_ie_i^2}{S_{yy}}.
\]
一元含截距回归中，\(\widehat y_i-\bar y=\widehat b(x_i-\bar x)\)，所以\(R^2=S_{xy}^2/(S_{xx}S_{yy})=r^2\)。教学数据中，\(600=360+240\)，故\(R^2=0.6\)。

这里的普通内积与前面有限概率空间的加权内积，分别来自数据的等权平方和与模型的概率加权。两处的最小误差，都由已有的正交结构确定。

\begin{exercise}
\mbox{}\par\nobreak

（1）对两个解释变量\(x,z\)，拟合\(\widehat y=a+bx+cz\)。由残差与\(\boldsymbol1,\boldsymbol x,\boldsymbol z\)正交写出三个方程。说明这三个向量线性无关时系数为何唯一。

（2）证明给全部\(y_i\)加同一常数，不改变回归斜率、残差和\(R^2\)。

（3）给每对记录指定正权重\(w_i\)，改为最小化\(\sum_iw_i(y_i-a-bx_i)^2\)。推导残差所满足的两个方程，并用加权内积说明它仍然是正交投影问题。
\end{exercise}

\FloatBarrier
\newpage
\section{随机过程初步*}
\label{sec:prob-processes}

甲、乙两人投篮，命中者继续，未中则换人。两人每次命中率分别为\(0.6,0.8\)，不受以往情况影响；第一次由谁投篮各以\(1/2\)决定。随着投篮继续，每次由甲投篮的概率怎样变化？最终会趋于稳定吗？

以投篮者为当前状态，下一次由谁投篮的概率便由当前状态确定，无须保留更早的投篮记录。下面将其抽象为两状态模型，用条件概率、线性变换与数列递推研究。

\subsection{状态与转移}

设\(S_n\)表示第\(n\)步的状态，只能为\(A,B\)，\(n=0,1,\ldots\)。已知当前状态后，下一步的概率不再依赖更早的状态，并且各步使用同样的转移规律：
\[
 P(S_{n+1}=B\mid S_n=A)=\alpha,\qquad
 P(S_{n+1}=A\mid S_n=B)=\beta.
\]
这给出一个两状态、时间齐次的{\heiti Markov链}模型。保持原状态的概率分别为\(1-\alpha,1-\beta\)。

在投篮问题中，令\(A,B\)分别表示由甲、乙投篮，则\(\alpha=0.4\)、\(\beta=0.2\)。以第一次投篮前为第0步，第\(n\)步的状态便表示第\(n+1\)次的投篮者。

\begin{figure}[!ht]
\centering
\includegraphics[width=.58\linewidth]{probability-markov}
\caption{两状态转移模型；箭头表示下一步的条件概率。}
\end{figure}

记\(u_n=P(S_n=A)\)。按当前状态使用全概率公式，得到
\begin{equation}\label{eq:prob-two-state-recurrence}
 u_{n+1}=(1-\alpha)u_n+\beta(1-u_n)
 =\beta+(1-\alpha-\beta)u_n.
\end{equation}
状态本身随机，\(u_n\)则是按给定初始分布与模型计算出的数。

若以列向量记录分布，令
\[
 \boldsymbol p_n=\begin{pmatrix}u_n\\1-u_n\end{pmatrix},\qquad
 M=\begin{pmatrix}1-\alpha&\beta\\\alpha&1-\beta\end{pmatrix},
\]
于是\(\boldsymbol p_{n+1}=M\boldsymbol p_n\)，\(\boldsymbol p_n=M^n\boldsymbol p_0\)。各列元素非负、和为1，故更新后仍为概率向量。

\subsection{长期分布}

使分布一步更新后不变的概率向量称为{\heiti 不变分布}。若\(\alpha+\beta>0\)，解\(M\boldsymbol\pi=\boldsymbol\pi\)，得到
\begin{equation}\label{eq:prob-stationary}
 \boldsymbol\pi=\begin{pmatrix}\beta/(\alpha+\beta)\\\alpha/(\alpha+\beta)\end{pmatrix}.
\end{equation}
记\(u_*=\beta/(\alpha+\beta)\)、\(r=1-\alpha-\beta\)。从递推式中减去固定点，
\[
 u_{n+1}-u_*=r(u_n-u_*),\qquad
 u_n=u_*+r^n(u_0-u_*).
\]
当\(0<\alpha+\beta<2\)时，\(|r|<1\)，故任意初始分布都趋于上述不变分布。

不变与收敛是不同任务。取\(\alpha=\beta=1\)，状态每步交替，仍有不变分布\((1/2,1/2)^{\mathsf T}\)；但从确定状态开始，分布在两个确定状态之间往返而不收敛。若初始分布恰为不变分布，则各时刻分布保持不变。

前\(n\)步处于状态\(A\)的总次数为\(N_n=\sum_{j=0}^{n-1}I_j\)，其中\(I_j=1_{\{S_j=A\}}\)。由期望的线性性与递推通项，无须各步独立便有
\[
 EN_n=\sum_{j=0}^{n-1}u_j
 =nu_*+(u_0-u_*)\frac{1-r^n}{1-r}
 \quad(\alpha+\beta>0).
\]
\begin{example}[据2023年新课标I卷改编]
沿用本节开头的投篮模型。

（1）求第二次由乙投篮的概率。

（2）求第\(i\)次由甲投篮的概率，并研究其极限。

（3）求前\(n\)次投篮中甲投篮次数的期望。说明这里为什么不能直接使用同一参数的二项分布。
\end{example}

\begin{exercise}
\mbox{}\par\nobreak

（1）取\(\alpha=0.3,\beta=0.2\)，初始状态为\(A\)。求\(u_n\)、不变分布及前\(n\)步处于\(A\)的次数期望。

（2）分别取\(\alpha=\beta=0\)与\(\alpha=\beta=1\)，比较不变分布的唯一性与分布收敛性。再说明\(\alpha=0,0<\beta<1\)时模型的长期行为。
\end{exercise}

从可能结果建立概率分配，从数值记录研究分布，再由抽样把模型与数据联系起来，本章的问题依次得到各自的表达。在最后这个模型中，条件分解、线性变换、期望与极限又共同工作。它们所提供的依据，已经足以完成当前的研究。

